\IfFileExists{stacks-project.cls}{%
\documentclass{stacks-project}
}{%
\documentclass{amsart}
}

% The following AMS packages are automatically loaded with
% the amsart documentclass:
%\usepackage{amsmath}
%\usepackage{amssymb}
%\usepackage{amsthm}

\usepackage{amssymb}

% For dealing with references we use the comment environment
\usepackage{verbatim}
\newenvironment{reference}{\comment}{\endcomment}
%\newenvironment{reference}{}{}
\newenvironment{slogan}{\comment}{\endcomment}
\newenvironment{history}{\comment}{\endcomment}

% For commutative diagrams you can use
% \usepackage{amscd}
\usepackage[all]{xy}

% We use 2cell for 2-commutative diagrams.
\xyoption{2cell}
\UseAllTwocells

% To put source file link in headers.
% Change "template.tex" to "this_filename.tex"
% \usepackage{fancyhdr}
% \pagestyle{fancy}
% \lhead{}
% \chead{}
% \rhead{Source file: \url{template.tex}}
% \lfoot{}
% \cfoot{\thepage}
% \rfoot{}
% \renewcommand{\headrulewidth}{0pt}
% \renewcommand{\footrulewidth}{0pt}
% \renewcommand{\headheight}{12pt}

\usepackage{multicol}

% For cross-file-references
\usepackage{xr-hyper}

% Package for hypertext links:
\usepackage{hyperref}

% For any local file, say "hello.tex" you want to link to please
% use \externaldocument[hello-]{hello}
\externaldocument[introduction-]{introduction}
\externaldocument[representationtheory-]{representationtheory}
\externaldocument[representations-compact-]{representations-compact}
\externaldocument[liegroups-general-]{liegroups-general}
\externaldocument[liestructure-]{liestructure} 
\externaldocument[algebraicgroups-]{algebraicgroups}
\externaldocument[reductiveforms-]{reductiveforms}
\externaldocument[vermamodules-]{vermamodules}
\externaldocument[gKmodules-]{gKmodules}
\externaldocument[asymptotics-]{asymptotics}
\externaldocument[plancherel-]{plancherel}
\externaldocument[discreteseries-]{discreteseries}
\externaldocument[galoiscohomology-]{galoiscohomology}
\externaldocument[automorphicspace-]{automorphicspace}
%\externaldocument[harmonicanalysis-]{harmonicanalysis} 
%\externaldocument[automorphicforms-]{automorphicforms}
%\externaldocument[periods-]{periods}
%\externaldocument[traceformulalocal-]{traceformulalocal}
%\externaldocument[traceformulaglobal-]{traceformulaglobal}
%\externaldocument[arithmetic-]{arithmetic}
%\externaldocument[geometric-]{geometric}
\externaldocument[fdl-]{fdl}
\externaldocument[index-]{index}

% Theorem environments.
%
\theoremstyle{plain}
\newtheorem{theorem}[subsection]{Theorem}
\newtheorem{proposition}[subsection]{Proposition}
\newtheorem{lemma}[subsection]{Lemma}

\theoremstyle{definition}
\newtheorem{definition}[subsection]{Definition}
\newtheorem{example}[subsection]{Example}
\newtheorem{exercise}[subsection]{Exercise}
\newtheorem{situation}[subsection]{Situation}

\theoremstyle{remark}
\newtheorem{remark}[subsection]{Remark}
\newtheorem{remarks}[subsection]{Remarks}

\numberwithin{equation}{subsection}

% Macros
%
\def\lim{\mathop{\rm lim}\nolimits}
\def\colim{\mathop{\rm colim}\nolimits}
\def\Spec{\mathop{\rm Spec}}
\def\Hom{\mathop{\rm Hom}\nolimits}
\def\SheafHom{\mathop{\mathcal{H}\!{\it om}}\nolimits}
\def\SheafExt{\mathop{\mathcal{E}\!{\it xt}}\nolimits}
\def\Sch{\textit{Sch}}
\def\Mor{\mathop{\rm Mor}\nolimits}
\def\Ob{\mathop{\rm Ob}\nolimits}
\def\Sh{\mathop{\textit{Sh}}\nolimits}
\def\NL{\mathop{N\!L}\nolimits}
\def\proetale{{pro\text{-}\acute{e}tale}}
\def\etale{{\acute{e}tale}}
\def\QCoh{\textit{QCoh}}
\def\Ker{\text{Ker}}
\def\Im{\text{Im}}
\def\Coker{\text{Coker}}
\def\Coim{\text{Coim}}

\def\eqref #1{(\ref{#1})}


% OK, start here.
%
\begin{document}

\title{Representations of reductive groups over local fields}


\maketitle

\phantomsection
\label{section-phantom}


\tableofcontents

[This chapter, especially the theory of asymptotics, is under construction and has not been proofread. Some statements may be slightly imprecise.]

In this chapter, we discuss representations of a group of the form $G(F)$, when $F$ is a local (locally compact) field, and $G$ is a reductive algebraic group over $F$. We treat the Archimedean and non-Archimedean cases in parallel, highlighting similarities. For economy of language, such a group will be called a ``real or $p$-adic reductive group'' --- the $p$-adic case including non-Archimedean local fields in equal characteristic, $F=\mathbb F_q((t))$. Everything in this chapter also applies to finite central extensions of reductive groups of the form $G(F)$, like the \emph{metaplectic group}, which are not necessarily algebraic; however, the notation is mostly adapted to the algebraic case.  When it is clear from the context, the group $G(F)$ will simply be denoted by $G$. If the word ``reductive'' is omitted, a ``real group'' will be a Lie group, and a ``$p$-adic group'' will be a $p$-adic analytic group (although most statements will be true for arbitrary totally disconnected, locally compact groups, in this case).

\begin{remark}
 \label{remark-complexification-Lie-algebra}
The ``real'' case includes the case when $F=\mathbb C$ --- notice that the complex structure plays no role in setting up the representation-theoretic problems, and we can think instead of $G(\mathbb C)$ as $\text{Res}_{\mathbb C/\mathbb R} G(\mathbb R)$. A common misunderstanding, when $G=G(\mathbb C)$ is a complex group, and $(\pi,V)$ is a smooth complex representation of $G$, is that the (complex) Lie algebra $\mathfrak g$ acts by complex-linear endomorphisms on $V$; it does not! Instead, $G$ should be treated as a real Lie group; for any smooth, complex representation of a real Lie group, we have an action of the complexified Lie algebra $\mathfrak g\otimes_{\mathbb R} \mathbb C$ on $V$ by complex-linear automorphisms. Since $\mathfrak g\otimes_{\mathbb R} \mathbb C \simeq \mathfrak g \oplus \mathfrak g$, what in the literature of representations of real reductive groups is a $\mathfrak g$-module, when one specializes to complex reductive groups it becomes a $\mathfrak g$-bimodule. 

To avoid confusion, in this chapter we will replace every complex group by its restriction of scalars over $\mathbb R$. Also, as is customary, we will often write $\mathfrak g$ (but will say so) for the \emph{complexification} of the Lie algebra $\mathfrak g_{\mathbb R}$ of any real Lie group $G$, i.e., $\mathfrak g := \mathfrak g_{\mathbb R}\otimes_{\mathbb R} \mathbb C$. 
\end{remark}


\section{Various categories of representations}
\label{section-categories-representations}

\subsection{Smooth and SF-representations}
\label{subsection-smooth-SF-representations}


The notion of a continuous, in particular of a Banach representation of a topological group was introduced in Definition \ref{representationtheory-definition-representation}. We also introduced an $F$-representation (or Fr\'echet representation of moderate growth) in Definition \ref{representationtheory-definition-Frepresentation}, which is a Fr\'echet representation that is a countable limit of Banach representations. 

\begin{definition}
 \label{definition-smooth-representation}
A {\it smooth vector} in a Banach or Fr\'echet representation $(\pi, V)$ of a real or $p$-adic group, resp.\ an {\it analytic vector}, in the real case, is a vector $v\in V$ such that the action map $G\ni g \mapsto \pi(g) v \in V$ is smooth (resp., analytic).\footnote{In the $p$-adic case, ``smooth'' means locally constant, so the definition is equivalent to requiring that $v$ have an open stabilizer.} In particular, in the real case, for a smooth vector $v$ and any element $D$ of the universal enveloping algebra $U(\mathfrak g_{\mathbb C})$, the element $\pi(D) v$ is defined.

The {\it space of smooth vectors} of a representation $(\pi, V)$ is denoted by $V^\infty$, and considered as a topological space, in the $p$-adic case with the direct limit topology over the subspaces $V^J$ as $J$ varies over open compact subgroups, and in the real case with the topology of convergence of all $\pi(D)v$, where $D$ ranges over elements of the universal enveloping algebra $U(\mathfrak g_{\mathbb C})$. 
 
A {\it smooth representation} $(\pi, V)$ is a representation such that $V = V^\infty$ as topological vector spaces. 

An {\it SF-representation}, or {\it smooth representation of moderate growth} of a real or $p$-adic group $G$ is a smooth F-representation. 
\end{definition}

\begin{lemma}
 \label{lemma-smooth-dense}
If $V$ is a Fr\'echet representation of a Lie or $p$-adic group $G$, the subspace $V^\infty$ of smooth vectors is dense.
\end{lemma}

\begin{proof}
 By Proposition \ref{representationtheory-proposition-integral-lcs}, the algebra $M_c^\infty(G)$ of smooth, compactly supported measures (= smooth, compactly supported functions times a Haar measure) acts on $V$. The image of the action is clearly in $V^\infty$, and by an approximation of the identity, one sees that the image is dense. 
\end{proof}

A much stronger, and important, statement is true: the {\it Dixmier--Malliavin theorem} states that the image of the action of $M_c^\infty(G)$ is all of $V^\infty$:

\begin{theorem}
 \label{theorem-Dixmier-Malliavin}
Let $V$ be a Fr\'echet representation of a Lie group or $p$-adic group $G$. The action map 
$$ M_c^\infty(G)\otimes V \to V^\infty$$
is surjective.
\end{theorem}

Notice that the tensor product here is not completed! The theorem means that every smooth vector can be written as a finite linear combinations of smooth, compactly supported measures acting on other vectors. (Also, without loss of generality, one might assume that $V=V^\infty$, if desired.)

\begin{proof}
The $p$-adic case is trivial, since every $J$-invariant vector (where $J$ is an open compact subgroup) is fixed by the action of $e_J=$ the probability Haar measure on $J$. The real case is the theorem of Dixmier--Malliavin, see \cite{Dixmier-Malliavin}, or \cite{Casselman-Dixmier-Malliavin}.
\end{proof}



\begin{remark}
\label{remark-smooth-not-SF}
Outside of the realm of F-representations (Fr\'echet representations of moderate growth), the notion of smooth representation leads to counterintuitive examples, e.g., the space of distributions on a Lie group $G$ is a smooth representation. We will only be considering smooth Fr\'echet representations of moderate growth from now on.
\end{remark}

\begin{lemma}
 \label{lemma-smooth-vectors-Frepresentation}
If $V$ is an F-representation of a real group, then $V^\infty$ is an SF-representation.
\end{lemma}

\begin{proof}
First of all, notice that the topology on $V^\infty$ is also given by a countable set of $G$-continuous seminorms: If $\rho_n$ is a sequence of $G$-continuous seminorms on $V$, defining its topology, and we fix, for every $d\ge 0$, a basis $(D_{d,i})_i$ of the $d$-th filtered part of the universal enveloping algebra $U(\mathfrak g_{\mathbb C})$, then the seminorms $\rho_{d,n}(v) = \max_i \rho_n (D_{d,i}v)$ define the topology on $V^\infty$ as $n$ and $d$ vary, and are $G$-continuous, because $\rho_{d,n}(gv) = \max_i \rho_n (g \cdot \text{Ad}(g^{-1})(D_{d,i}) v) \ll \rho_{d,n}(v)$ (locally uniformly in $G$), since the adjoint representation preserves the filtration. 


The content of the lemma, then, is that the topological vector space $V^\infty$ is complete. One shows\footnote{\label{footnote-smooth}We will be abusing notation and writing $C^\infty(G, V)$ for $C(G, V)^\infty$, the space of smooth vectors in the space of continuous functions from $G$ to $V$. For example, in the $p$-adic case, these are not just locally constant functions but \emph{uniformly} locally constant, in the sense that they are fixed by a compact open subgroup.} that the action map $g\mapsto \pi(g) v$ gives rise to a morphism $V^\infty \to C^\infty(G, V)$, where $G$ acts by the right regular representation on $C^\infty(G, V)$, and that this is an isomorphism onto the closed subspace $C^\infty(G, V)^G$ of functions that are invariant under the simultaneous action: $g\cdot f (x) := \pi(g) f(g^{-1} x)$. 
\end{proof}

\subsection{Unitary representations}
\label{subsection-unitary-representations}

Unitary representations have been introduced in \S \ref{representationtheory-section-unitary-representations}. Their Plancherel decomposition was discussed in \S \ref{representationtheory-section-Plancherel}. Here, we will just add the uniqueness of the Plancherel decomposition, for reductive real or $p$-adic groups. [LATER]


\subsection{$(\mathfrak g, K)$-modules}
\label{subsection-gK-modules}

Topological representations of Lie groups do not form an abelian category. This is sometimes cumbersome; to make the theory more algebraic, we sometimes work with $(\mathfrak g,K)$-modules. 

\begin{definition}
 \label{definition-g-K-module}
Let $\mathfrak g$ be a complex Lie algebra, and $H$ a Lie group, with an embedding $\mathfrak h_{\mathbb C}\hookrightarrow \mathfrak g$, and a representation $\text{Ad}:H \to \text{GL}(\mathfrak g)$, extending the adjoint action on $\mathfrak h_{\mathbb C}$, whose differential coincides with the adjoint action of $\mathfrak h\subset\mathfrak g$. (For example, $\mathfrak g$ is the complexified Lie algebra of a Lie group containing $H$.) A {\it $(\mathfrak g, H)$-module} is a vector space $V$ with actions of both $\mathfrak g$ and $H$, such that:
\begin{enumerate}
 \item the action of $H$ is locally finite;
 \item the differential of the action of $H$ coincides with the action of $\mathfrak h$, considered as a subalgebra of $\mathfrak g$;
 \item $h\cdot X\cdot h^{-1}\cdot v = \text{Ad}(h)(X)\cdot v$, for all $h\in H$, $X\in \mathfrak g$, $v\in V$.
\end{enumerate}
\end{definition}

\begin{remark}
 \label{remark-gK-algebraic}
This notion is most often (but not exclusively!) used when $H=K$ is a maximal compact subgroup of a Lie group $G$ (with complexified Lie algebra $\mathfrak g$). In that case, by Theorems \ref{liegroups-general-theorem-Weyl}, \ref{reductiveforms-theorem-functor-compactgroups}, $K=\mathbf K(\mathbb R)$ for an anisotropic reductive algebraic group $\mathbf K$, and there is no difference between locally finite $K$-modules and locally algebraic $\mathbf K$-modules. Thus, in that case, the notion of a $(\mathfrak g, K)$-module can be defined in a completely algebraic way, using the complex pair $(\mathfrak g, \mathbf K)$, and without any explicit reference to the real structure. 
\end{remark}

\begin{lemma}
 \label{lemma-gK-of-representation}
Let $(\pi, V)$ be a representation of a Lie group $G$, and $H\subset G$ a subgroup. The subspace $V_{H\text{-fin}}$ of $H$-finite vectors is stable under the action of $\mathfrak g_{\mathbb C}$.
\end{lemma}

\begin{proof}
 For every $v\in V_{H\text{-fin}}$, the image of the action map $\mathfrak g \otimes \text{span}(Hv) \to V$ is finite-dimensional, and contains the element $h\cdot X\cdot v$ for all $X\in \mathfrak g$ and $h\in H$, since $h\cdot X\cdot v = \text{Ad}(h)(X)\cdot h\cdot v$.
\end{proof}


Recall also from Theorem \ref{representations-compact-theorem-PeterWeyl-general} that if $H=K$ is compact, and the representation is Fr\'echet, the space of $K$-finite vectors is dense. 



\begin{definition}
 \label{definition-gK-of-representation}
Let $G$ be a reductive Lie group, and $K\subset G$ a maximal subgroup; use $\mathfrak g$ to denote the complexified Lie algebra of $G$. The {\it $(\mathfrak g, K)$-module of a Fr\'echet representation} $(\pi, V)$ of $G$ is the  $(\mathfrak g, K)$-module $V^\infty_{K\text{-fin}}$ of $K$-finite smooth vectors in $V$. 

Two representations $V_1$, $V_2$ are said to be {\it infinitesimally equivalent} if their $(\mathfrak g, K)$-modules are isomorphic.
\end{definition}


\begin{remark}
 \label{remark-infinitesimal-equivalence}
Infinitesimal equivalence captures more of the essence of representation theory than isomorphisms of representations. For example, all Banach representations $L^p(\mathbb R^\times)$ ($p\ge 1$) of the group $\mathbb R^\times$ are infinitesimally equivalent, although they are not isomorphic as topological vector spaces. On the other hand, the ``globalization'' theorem of Casselman and Wallach \cite{Casselman-canonicalextensions, Wallach-RR2, Bernstein-Kroetz}) says that any finitely generated, \emph{admissible} (see Definition \ref{definition-admissible}) $(\mathfrak g, K)$-module admits a unique ``globalization'' to a smooth Fr\'echet representation of moderate growth. The proof of this theorem relies on the subrepresentation theorem (see Theorem \ref{theorem-subrepresentation}), realizing irreducible $(\mathfrak g,K)$-modules as submodules of parabolically induced representations.
\end{remark}



\begin{lemma}
 \label{lemma-gK-dense}
If $V$ is a Fr\'echet representation of a reductive Lie group $G$, and $K\subset G$ a maximal subgroup, its $(\mathfrak g,K)$-module $V^\infty_{K\text{-fin}}$ is dense in $V$. In particular, if the $(\mathfrak g, K)$-module $V^\infty_{K\text{-fin}}$ is irreducible, so is $V$. 
\end{lemma}

\begin{proof}
 This follows from Lemma \ref{lemma-smooth-dense} and Proposition \ref{representations-compact-proposition-denseinclusions}. 
\end{proof}

The converse is true in the category of \emph{admissible} representations (Theorem \ref{theorem-irreducible-admissible}). 


\subsection{Admissibility}
\label{subsection-admissibility}

Let $G$ be a real or $p$-adic reductive group. In the real case, let $K\subset G$ be a maximal compact subgroup, and denote by $\mathfrak g$ the complexification of the Lie algebra.

\begin{definition}
 \label{definition-admissible}
A $(\mathfrak g, K)$-module $(\pi, V)$ (in the real case), or a smooth $K$-module $V$ (in the $p$-adic case) is called {\it admissible} if  all irreducible representations of $K$ appear with finite multiplicity, i.e., $\dim\Hom_K(\tau, V)<\infty$ for every irreducible representation $\tau$ of $K$. 

A (topological) representation $(\pi, V)$ of a real or $p$-adic reductive group $G$ is admissible if the $(\mathfrak g, K)$-module (resp.\ $K$-module, in the $p$-adic case) $V^\infty_{K\text{-fin}}$ is admissible. Here, $K$ is any maximal compact subgroup of $G$, in the real case, and any compact open subgroup of $G$, in the $p$-adic case.
\end{definition}


\begin{remark}
 \label{remark-admissible-indep-K}
The property of being admissible, for a representation of $G$, does not depend on the choice of $K$; indeed, in the real reductive case, all Cartan subgroups are conjugate, by Theorem \ref{reductiveforms-theorem-Cartan-involution-exists}. In the $p$-adic case, the independence follows from the lemma below.
\end{remark}

\begin{lemma}
 \label{lemma-admissible-padic}
In the $p$-adic case, a representation $(\pi, V)$ is admissible if and only if, for every compact open $J\subset G$, we have $\dim V^J <\infty$. 
\end{lemma}

\begin{proof}
 First of all, observe that $V^\infty_{K\text{-fin}} = V^\infty$ for every compact open $K\subset G$.  
 
 If a (smooth) irreducible representation $\tau$ of $K$ appears with infinite multiplicity, then, obviously, $\dim V^J = \infty$ for all $J$ with $\tau^J\ne 0$. 
 
 Vice versa, given $K$, for every open compact $J\subset K$, the set of (isomorphism classes of) irreducible representations $\tau$ of $K$ with $\tau^J\ne 0$ is finite. Indeed, to prove this claim, we can replace $J$ with the intersection of all its $K$-conjugates, which is still open and compact, but also normal. Then, if $\tau^J\ne 0$ and $\tau$ is irreducible, we have $\tau = \tau^J$, hence $\tau$ is an irreducible representation of the finite group $K/J$, and there are only finitely many such. Thus, admissibility according to Definition \ref{definition-admissible} implies that $V^J$ is finite-dimensional, for every $J$.  
\end{proof}

\begin{definition}
 \label{definition-contragredient}
The {\it contragredient} of a $(\mathfrak g,K)$-module $V$, in the real case, or a smooth $G$-representation $V$, in the $p$-adic case, is the $(\mathfrak g, K)$-module, resp.\ smooth $G$-representation $\tilde V:=(V^*)_{K\text{-fin}}$ of $K$-finite vectors in the linear dual of $V$. 
\end{definition}



\begin{lemma}
 \label{lemma-admissible-gK}
A $(\mathfrak g,K)$-module $V$, in the real case, or a smooth $G$-representation $V$, in the $p$-adic case, is admissible if and only if the natural embedding $V\hookrightarrow \tilde{\tilde V}$ is an isomorphism.

If $V$ is admissible and irreducible, any automorphism of $V$ (as a $(\mathfrak g,K)$-module, resp.\ as a $G$-representation) is scalar.
\end{lemma}

\begin{proof}
 The module is a direct sum over its $K$-types, $V = \bigoplus_\tau V_{K,\tau}$, where $V_{K,\tau}$ denotes the $\tau$-isotypic subspace, and $\tau$ ranges over all irreducible representations of $K$. The contragredient is then $\tilde V = \bigoplus_\tau V_{K,\tau}^*$, and the first claim follows from linear algebra. 
 
 Any automorphism preserves the isotypic spaces. If those are finite-dimensional, then the automorphism will have an eigenvector, and since this eigenvector generates the entire $(\mathfrak g, K)$-module, it immediately follows that the automorphism is scalar.
\end{proof}

The converse to Lemma \ref{lemma-gK-dense} holds, for admissible representations: 

\begin{theorem}
 \label{theorem-irreducible-admissible}
If $V$ is an irreducible admissible Fr\'echet representation of moderate growth of a reductive Lie group $G$, its $(\mathfrak g, K)$-module is irreducible, and all $K$-finite vectors are automatically analytic (in particular, smooth).
\end{theorem}

\begin{proof}
First of all, since $V^\infty_{K\text{-fin}}$ is dense (Lemma \ref{lemma-gK-dense}) in $V$, it is also dense in $V_{K\text{-fin}}$. For any $K$-type $\tau$, there is a measure $\mu_\tau$ on $K$ whose action on any Fr\'echet module is a projection onto the $\tau$-isotypic component. Therefore, the $\tau$-isotypic subspace $V^{\infty,\tau}$ is dense in $V^\tau$. But the former is finite-dimensional, therefore the two coincide, i.e., every $K$-finite vector is smooth.

Suppose that $V_0\subset V_{K\text{-fin}}$ is a nonzero $(\mathfrak g, K)$-submodule. We claim that the closure of $V_0$ is $G$-stable. This requires the ``big hammer'' of elliptic regularity to prove, so we only give a couple of steps, followed by references.

First, we notice that the action of the center $\mathfrak z(\mathfrak g)$ of the universal enveloping algebra of the (complexified) Lie algebra $\mathfrak g$ on $V_{K\text{-fin}}$ is locally finite: indeed, it preserves the finite-dimensional, $K$-isotypic subspaces. 

Elliptic regularity, now, implies that all vectors in $V_{K\text{-fin}}$ are analytic; see \cite[3.4.9]{Wallach-RR1}.\footnote{In Wallach's book, the argument is formulated for representations on Hilbert spaces, but it holds verbatim for Banach spaces, and hence for Fr\'echet representations of moderate growth. Note that a function $G\to V$, where $V$ is a Banach space, is (real) analytic iff it is \emph{weakly} analytic, i.e., iff its composition with any continuous functional $v^*: V\to \mathbb C$ is analytic.} And, the closure of a $\mathfrak g$-stable subspace of analytic vectors in $V$ is stable under the identity component of $G$: simply apply the exponential map $\mathfrak g_{\mathbb R} \to G$, whose image generates the identity component. 

Since $V_0$ is not only $\mathfrak g$-stable, but also $K$-stable, and $K$ meets all connected components of $G$, $V_0$ is $G$-stable. Since $V$ is irreducible, $V_0$ is dense. But, again, applying projectors to the $K$-types, this means that for any $K$-type $\tau$, the $\tau$-isotypic subspace $V_0^\tau$ is dense in $V^\tau$. Since these spaces are finite-dimensional, $V_0^\tau=V^\tau$ for all $\tau$, hence $V_0 = V_{K\text{-fin}}$. 
\end{proof}





\section{Schwartz and Harish--Chandra Schwartz spaces}
\label{section-Schwartz}


\subsection{Schwartz space defined by a scale function}
\label{subsection-Schwartz-scale}

We follow \cite[\S 2]{Bernstein-Kroetz}.

\begin{definition}
 \label{definition-scale}
A {\it scale} on a locally compact group $G$ is a function $s : G \to \mathbb R^+$ such that:
\begin{itemize}
 \item $s$ and $s^{-1}$ are locally bounded,
 \item $s$ is submultiplicative, i.e., $s(gh) \le s(g)s(h)$ for all $g, h \in G$.
\end{itemize}

A scale function $s'$ \emph{dominates} a scale function $s$, if there exist positive constants $C, N$ such that $s \le C s'^N$. They are \emph{equivalent} if each dominates the other. 


A {\it scale structure} on $G$ is an equivalence class of scale
functions.
\end{definition}

In other words, a scale function is the exponential of a radial function, Definition \ref{representationtheory-definition-radial-function}.


In \ref{representationtheory-definition-radial-natural} we saw the ``natural radial function'' $r_{\text{nat}}$ (and hence its exponential, the ``natural scale function'' $s_{\text{nat}}$, denoted $\Vert g\Vert$ there) of a compactly generated group; recall that $r_{\text{nat}}(g)$ counts how many times we need to multiply a compact generating neighborhood of the identity by itself in order to produce a set containing $g$. 


\begin{definition}
 \label{definition-Schwartz-space}
Let $G$ be a group equipped with a scale structure $[s]$ (Definition \ref{definition-scale}), with associated radial function $r=\log s$. The associated {\it Schwartz space} $\mathcal S_{[s]}(G)$ is the space of smooth vectors in the left and right regular F-representation on measures $f$ on $G$ which satisfy $f \cdot s^n \in L^1(G)$ for all $n\in \mathbb N$. 

 The {\it natural Schwarz space} $\mathcal S_{\text{nat}}(G)$ is the one defined by the class of natural scale functions. 
\end{definition}

Equivalently, the natural Schwartz space is the space of smooth vectors in the space of rapidly decaying measures of Definition \ref{representationtheory-definition-rapidly-decaying}. 


\begin{example}
 \label{example-natural-Schwartz-additive-multiplicative}
For the additive group $G=\mathbb G_a(F)$, we have $\mathcal S_{nat}(G)=$ the Schwartz space of smooth functions (times a Haar measure) which, together with their derivatives (in the real case), are of superexponential decay. On the other hand, for $G=\mathbb G_m(F)$, they coincide with smooth functions $f$ (times a Haar measure) such that $f(x)\cdot |x|^n$ is bounded for all $n\in \mathbb Z$, and similarly for all derivatives (in the real case).
\end{example}

\begin{remark}
 \label{remark-almost-smooth}
There is some clumsiness in trying to deal with the real and $p$-adic cases at the same time, which is due to the fact that the notion of ``smooth'' in the $p$-adic case is not quite analogous to that of ``smooth'' in the real case; for example, smooth vectors in an F-representation of a $p$-adic group do not produce Fr\'echet spaces. There is a notion of ``almost smooth'' vectors in the $p$-adic case, which is a better analogy to ``smooth'' in the real case, see \cite{Sakellaridis-beyond1}, but it is not very useful in practice. Because of the strong definition of smoothness (=local constancy), and because we tend to forget about the topology on spaces of smooth vectors of representations of $p$-adic groups, the ``rapid decay'' Schwartz spaces that we are defining here are not suitable for $p$-adic groups; in the next subsection, we will discuss algebraically defined Schwartz spaces using compactifications, where the definitions in the real and $p$-adic cases coincide, and produce compactly supported functions/measures in the $p$-adic case. [But, note for the future: Maybe we can expand the notion of SF-representation to the $p$-adic case, to include the LF-spaces of smooth vectors in an F-representation; or, include a full discussion of ``almost smooth'' vectors, for the sake of uniformity.]
\end{remark}



\begin{proposition}
 \label{proposition-SF-Schwartz-equivalence}
Let $G$ be a real Lie group. The categories of smooth Fr\'echet representations of moderate growth of $G$, and of nondegenerate continuous algebra representations of $\mathcal S_{\text{nat}}(G)$ on Fr\'echet spaces, are equivalent.
\end{proposition}

\begin{proof}
 If $(\pi, V)$ is any F-representation, the action of $G$ extends to a continuous representation of the algebra of rapidly decaying measures by Proposition \ref{representationtheory-proposition-rapiddecay-Banach}, in particular, to the natural Schwartz space. 

 A theorem of Dixmier and Malliavin \cite{Dixmier-Malliavin} states that, if $(\pi,V)$ is a smooth Fr\'echet representation of a real Lie group $G$, then the action map $M_c^\infty(G) \otimes V \to V$ is surjective. Hence, so is the map $\mathcal S_{\text{nat}}(G)\otimes V \to V$, i.e., $V$ is nondegenerate. 
 
 Vice versa, if $V$ is a nondegenerate continuous Fr\'echet $\mathcal S_{\text{nat}}(G)$-module, that is, it is nondegenerate and the action map $\mathcal S_{\text{nat}}(G) \times V\to V$ is continuous, this action extends to the \emph{projective tensor product} $\mathcal S_{\text{nat}}(G) \hat\otimes_\pi V \to V$, which is also a Fr\'echet space, and this gives a topological identification of $V$ as a quotient of the projective tensor product. Quotients of SF-representations are SF-representations, see \cite[Lemma 2.9 and Proposition 2.20]{Bernstein-Kroetz} for more details. 
\end{proof}





\subsection{Schwartz space of a semi-algebraic manifold}
\label{subsection-Schwartz-algebraic}

If $G$ denotes the points of a linear algebraic group over a local field, we also define another scale function, that depends on the algebraic structure. (The same definition can be given for finite covers thereof, by passing to the algebraic quotient.)

\begin{definition}
 \label{definition-algebraic-scale}
Let $G$ be a linear algebraic group, and fix a closed embedding $G\hookrightarrow \mathbb A^r$, with coordinates $x_1, \dots, x_r$. The corresponding {\it algebraic scale function} of $G(F)$, where $F$ is a local field, is 
$$s_{\text{alg}} (g) = \max_i |x_i(g)|.$$
\end{definition}

It is easy to prove that any two algebraic scale functions are equivalent.

\begin{lemma}
 \label{lemma-reductive-algebraic-natural}
If $G$ is a reductive group, the natural and algebraic scale functions on $G$ are equivalent.
\end{lemma}

\begin{proof}
 The statement is easily seen to be true for a torus. For a general reductive group, it reduces to the case of tori by the Cartan decomposition $G=K A^+ K$.
\end{proof}

This leads to a notion of ``algebraic Schwartz space'' according to Definition \ref{definition-Schwartz-space}, but in the $p$-adic case we would like a stricter definition that coincides with the space of compactly supported smooth measures. In this subsection, we will provide a uniform such for arbitrary real or $p$-adic (smooth) varieties (or semialgebraic spaces). 

[Definition of Schwartz space $\mathcal S(X)$ on a Nash manifold $X$ here. In the $p$-adic case, it coincides with $M_c^\infty(X)$. In particular, in the $p$-adic group case, $\mathcal S(G)=\mathcal H(G)=$ the Hecke algebra.]

\begin{proposition}
 \label{proposition-smooth-Schwartz-equivalence}
For both real and $p$-adic reductive groups, there is an equivalence of categories between SF-representations (in the real case), or smooth representations without topology (in the $p$-adic case), and nondegenerate $\mathcal S(G)$-modules.
\end{proposition} 

\begin{proof}
 In the real case, this is just Proposition \ref{proposition-SF-Schwartz-equivalence}, together with the equivalence of natural and algebraic scale structures, Lemma \ref{lemma-reductive-algebraic-natural}. 
 
 In the $p$-adic case, the proof is similar (but simpler). Notice that the analogous statement holds, more generally, for any locally compact, totally disconnected group.
\end{proof}

\subsection{Harish-Chandra Schwartz space}
\label{subsection-HC-Schwartz}

We follow \cite{Bernstein-Plancherel}. We start by defining a notion of radial function for a homogeneous space of a locally compact group; everything in this section applies to a finite union of homogeneous spaces, as well.


\begin{definition}
 \label{definition-radial-homogeneous}
Let $X$ be a homogeneous space for a locally compact group $G$. A {\it radial function} is a locally bounded function $r_X: X\to \mathbb R_+$, such that:
\begin{enumerate}
 \item for every $R \in \mathbb R_+$, the ``ball'' $B(R) = \{ x\in X| r_X(x)\le R\}$ is relatively compact in $X$;
 \item for any compact $\Omega\subset G$, there is a constant $C>0$ such that $|r_X(gx)-r_X(x)|<C$. 
\end{enumerate}
We say that $r_X$, $r_X'$ are two {\it equivalent radial functions} if there is a constant $C$ such that $C^{-1}(1+r_X)\le (1+r_X')\le C(1+r_X)$. 

We say that the space $X$ is of {\it polynomial growth} (with respect to a given equivalence class of radial functions) if there is a $d\ge 0$ such that, for one (any) compact neighborhood $\Omega$ of the identity in $G$, and some positive constant $C$, the ball $B(R)$ can be covered by $\le C(1+R^d)$ orbits of $\Omega$. 

The equivalence class of {\it natural radial functions} on $X$ is the equivalence class of the function $ r_X( x) = \inf\{r(g)| x_0\cdot g = x\}$, where $r$ is a natural radial function on $G$, and $x_0$ is some fixed point on $X$. 
\end{definition}

Now we assume that $X$ is a homogeneous real or $p$-adic manifold, of polynomial growth, under the action of a real or $p$-adic group $G$, or a finite union of such. [Fact, to be added: the space $X(F)$ of points of a spherical $G$-variety $X$ over a local field $F$, under the action of the group $G(F)$, are such.] 

\begin{definition}
 \label{definition-Harish-Chandra-Schwartz}
Let $X$ be a homogeneous real or $p$-adic manifold, under the action of a real or $p$-adic group $G$, of polynomial growth with respect to the natural scale. The {\it Harish-Chandra--Schwartz space} of $X$ is the space $\mathcal C(X)$ of smooth vectors in the Fr\'echet space of half-densities $f$ on $X$ with $f \in \lim_{\underset{d >0}{\leftarrow}} L^2(X, (1+r)^d)$, where $r$ is a natural scale function on $X$. 
\end{definition}

The notation $L^2(X, (1+r)^d)$ stands for the Hilbert space of half-densities $f$ with norm equal to the square root of $\int_X |f|^2 (1+r)^d$. 

\begin{remark}
 \label{remark-HCS-nuclear}
The space $\mathcal C(X)$ is a nuclear Fr\'echet space, in the real case, and a countable direct limit over the nuclear Fr\'echet spaces of $J$-invariants, as $J$ ranges over a basis of open compact subgroups, in the $p$-adic case.
\end{remark}



\begin{remark}
 \label{remark-half-densities}
If $X$ has an invariant measure $dx$, or, more generally, a positive $G$-eigenmeasure with (positive) $G$-eigencharacter $\eta$, one can think of half-densities as functions, by dividing by $(dx)^\frac{1}{2}$, but the action of $G$ on those functions is twisted by the square root of $\eta$, that is: 
\begin{equation}
 \label{equation-twisted-action}
 (g\cdot \Phi)(x) = \eta^\frac{1}{2}(g) \Phi(x\cdot g).
\end{equation}

In other words, if $\mathcal F(X)$ denotes functions and $\mathcal D(X)$ denotes half-densities, division by $(dx)^\frac{1}{2}$ defines an equivariant isomorphism $\mathcal D(X) \xrightarrow\sim \mathcal F(X) \otimes \eta^\frac{1}{2}$. 

For example, consider the pre-flag variety $X=U\backslash G_0$, where $U\subset P\subset G_0$ is the unipotent radical of a parabolic subgroup. Considering it as a homogeneous space for the product $G=L\times G_0$, where $L$ is the Levi quotient of $P$, it possesses $G$-eigenmeasure, which is invariant under $G_0$, but $\delta_P$-equivariant under $L$, where $\delta_P$ is the modular character of $P$. Thus, half-densities on $X$ can be identified (after a choice of such a measure, unique up to scalar), with functions on $X$, with the action of $L$ on the latter twisted by $\delta_P$.
\end{remark}


\begin{definition}
 \label{definition-tempered}
The space of {\it tempered half-densities} on $X$ is the dual of the topological vector space $\mathcal C(X)$. If a $G$-eigenmeasure $dx$ on $X$ is chosen (always to be taken $G$-invariant, if possible), the dual of the space $(dx)^{-\frac{1}{2}} \mathcal C(X)$ of Harish-Chandra--Schwartz functions is the space of {\it tempered measures}, and the dual of the space $(dx)^{\frac{1}{2}} \mathcal C(X)$ of Harish-Chandra--Schwartz measures is the space of {\it tempered generalized functions}. 

The space of {\it tempered smooth half-densities} (and, correspondingly, functions or measures in the presence of an eigenmeasure) is the space of smooth vectors in the contragredient of the F-representation $\bigcap_d L^2(X, (1+r)^d)$ (Definition \ref{representationtheory-definition-contragredient-F}), that is, in the direct limit of Hilbert spaces $\lim_{\underset{d >0}{\to}} L^2(X, (1+r)^d)$.
\end{definition}

Here is the main result of \cite{Bernstein-Plancherel}:

\begin{theorem}
 \label{theorem-Bernstein-Plancherel}
The inclusion $\mathcal C(X)\hookrightarrow L^2(X)$ is {\it fine}; that is, for any morphism from $L^2(X)$ to a direct integral $H=\int H_z \mu(z)$ of Hilbert spaces, the composition $\mathcal C(X)\to H$ is pointwise defined (Definition \ref{representationtheory-definition-pointwise-defined}). 
\end{theorem}


\begin{proof}
 This is \cite[Theorem 3.2]{Bernstein-Plancherel}, applied to the setting of \cite[\S 3.5, 3.7]{Bernstein-Plancherel}.
\end{proof}

\begin{definition}
 \label{definition-tempered-representation}
An admissible smooth representation $\pi$ of a real or $p$-adic reductive group is called {\it tempered} if its matrix coefficients are tempered, i.e., have image in the space $C^\infty_{\text{temp}}(G)$ of smooth, tempered functions (Definition \ref{definition-tempered}).

More generally, if $X$ is a homogeneous $G$-space of polynomial growth, a morphism $\ell:\pi\to C^\infty(X)$ is called {\it tempered} if the image lies in $C^\infty_{\text{temp}}(X)$.
\end{definition}









\section{Asymptotics}
\label{section-asymptotics}

\subsection{General setup}
\label{subsection-asymptotics-setup}

When $X = H\backslash G$ is a homogeneous $G$-space, and $\pi$ a smooth representation of $G$, a morphism $m:\pi\to C^\infty(X)$ is sometimes called a \emph{generalized matrix coefficient}; the reason is that any such morphism is equivalent (by Frobenius reciprocity) to an $H$-invariant functional $\ell$, so $m(v)(x) = \left< \pi(g) v, \ell \right>$ is a ``matrix coefficient'', where the covector $\ell$ is allowed to be non-smooth. In this section, we compare generalized matrix coefficients of certain representations of $G$ on a spherical variety $X$, with generalized matrix coefficients on the boundary degenerations. 

There are similarities, but also differences, between the real and $p$-adic cases. The main difference, in the real case, is that we need to restrict to admissible modules. (A general theory of asymptotics for smooth representations would be very desirable, but has not yet been developed! The naive translation of statements from the $p$-adic to the real case does not hold, in general.)

For the remainder of this section, $G$ is a real or $p$-adic reductive group, and $K$ is a maximal compact subgroup, if $G$ is real. We compare generalized matrix coefficients on $X$ and $X_\Theta$ by choosing some reasonable (but noncanonical) identification of the spaces ``close to infinity''. We write $X=X(F)$, etc.

\begin{definition}
 \label{definition-approximate-exponential}
Let $Z$ be the closure of a $G$-orbit in a smooth toroidal embedding $\bar X$ of $X$. An {\it approximate exponential map} is an analytic map $\phi: U_Z\to \bar X,$ where $U_Z$ is a neighborhood of $Z$ in the $F$-points of the normal bundle $N_Z \bar X$,  with the property that the partial differential of $\phi$ induces the identity on $N_Z \bar X$, and $\phi$ maps the intersection of every $G$-orbit with $U_Z$ to the corresponding $G$-orbit on $\bar X$. The {\it exponential bundle} $\text{Exp}_Z \bar X$ over $Z$ is the fiber bundle of germs, over $Z$, of approximate exponential maps, i.e., approximate exponential maps defined in some neighborhood of $Z$, modulo the equivalence relation of being equal in a smaller neighborhood. 
\end{definition}

Note that $\text{Exp}_Z \bar X$ is a torsor for the group bundle $\text{Exp}_Z N_Z\bar X$ of germs of approximate exponential maps from the normal bundle to itself (defined the same way). 

\begin{proposition}
 \label{proposition-exponential-p-adic}
Assume that $F$ is non-Archimedean. Using the notation of Definition \ref{definition-approximate-exponential}, let $\phi: U_Z \to \bar X(F)$ be an approximate exponential map for some orbit closure $Z\subset \bar X$.  
Then, given an open compact subgroup $J\subset G$, there is a $J$-invariant neighborhood $U_Z'\subset U_Z$ of $Z$, with $J$-invariant image $U_X'\subset \bar X(F)$, such that $\phi$ descends to a bijection: $U_Z'/J\to U_X'/J$. Moreover, any two approximate exponential maps descend to the same bijection, if the neighborhood $U_Z'$ is taken sufficiently small.
\end{proposition}

\begin{proof}
 See \cite[Proposition 4.3.1]{Sakellaridis-Venkatesh}. The reader is encouraged to check it directly in the baby case of $\bar X=\mathbb A^1$, $Z=\{0\}$, $G=\mathbb G_m$.
\end{proof}

Now, the normal bundle to $Z$ contains some open $G$-orbit, which we have called the boundary degeneration; let's denote it by $X_\Theta$. This proposition implies that, for any $J$-invariant functions $f$, $f_\Theta$ on $X$ and $X_\Theta$, respectively, there is a well-defined notion of the functions being asymptotically equal:

\begin{definition}
 \label{definition-asymptotically-equal-padic}
Assume that $F$ is non-Archimedean. Let $X$ be a spherical variety, and $X_\Theta$ an asymptotic cone thereof, obtained as the open $G$-orbit in the normal bundle of some orbit $Z$ in a toroidal embedding. If $f\in C^\infty(X)$, $f_\Theta\in C^\infty(X_\Theta)$, we say that $f$ is {\it asymptotically equal} to $f_\Theta$, written $f\sim f_\Theta$, if there is an approximate exponential map $\phi: U_Z\to \bar X(F)$ (Definition \ref{definition-approximate-exponential}), where $U_Z$ is a neighborhood of $Z$, such that, after possibly replacing $U_Z$ by a smaller neighborhood, $\phi^* f|_{U_Z} = f_\Theta|_{U_Z}$.
\end{definition}

Notice that, by Proposition \ref{proposition-exponential-p-adic}, this notion does not depend on the choice of approximate exponential, i.e., if the statement is true for one such map, then it is true for all.

In the real case, things are finer, since smooth functions are not locally constant. Therefore, any such attempt to identify $f$ and $f_\Theta$ will depend on the choice of approximate exponential. Instead of looking at arbitrary smooth functions, here, we will restrict our attention to ``functions that look like generalized characters'' (of the tori $A_\Theta$) at infinity---we will call such functions ``asymptotically finite''. The following baby example captures the essence of such functions:

\begin{example}
 \label{example-real-asymptotics-baby}
Let $\bar X = \mathbb A^1 \supset X=\mathbb A^1\smallsetminus\{0\}$, over $F=\mathbb R$. Let $Z=\{0\}$; then, $N_Z \bar X = \mathbb A^1$. Here, we want to think of $\bar X$ simply as a variety (without a group action), while $N_Z \bar X$ has a $\mathbb G_m$-action. Any analytic map $\phi: U_Z\to \mathbb R$, where $U_Z$ is a neighborhood of zero, fixing zero and inducing the identity on its tangent space, is an asymptotic exponential. Explicitly, such a $\phi$ is given by a power series of the form $\phi(x) = x + \sum_{n=2}^\infty a_n x^n$, convergent within some radius. 

An ``asymptotically finite'' function $f$ on $X$ is a function with the property that $\phi^* f = \sum_\lambda f_\lambda \cdot h_\lambda$, a finite sum indexed by characters of the multiplicative group, where $f_\lambda$ is a generalized $\mathbb G_m$-eigenfunction with generalized eigencharacter $\lambda$, and $h_\lambda \in C^\infty(U_Z)$. The reader should check [exercise!] that this notion does not depend on the choice of approximate exponential $\phi$.
\end{example}

\begin{definition}
 \label{definition-asymptotically-finite}
 Let $F$ be  real or non-Archimedean, and let $X$ be a spherical variety over $F$. An \emph{asymptotically finite} function on $X$ is a smooth function $f$ with the property that, for some toroidal compactification $\bar X$, in a neighborhood of any point $z\in \bar X$ (belonging to a $G$-orbit $Z$ whose normal bundle is the boundary degeneration $X_Z$), and for any approximate exponential $\phi$ defined in a neighborhood $U$ of $z$, the function $\phi^*f$, restricted to a neighborhood $U'\subset U$ of $z$, is equal to 
 \begin{equation}
  \label{equation-asymptotically-finite}
  \sum_\lambda f_\lambda \cdot h_\lambda,
 \end{equation}
 a finite sum indexed by characters of $A_Z$, where $f_\lambda$ is a generalized $A_Z$-eigenfunction with generalized eigencharacter $\lambda$, and $h_\lambda \in C^\infty(U')$.
  
 We let $\text{Fin}_Z(\bar X)$ denote the bundle of germs, over a $G$-orbit $Z$, of asymptotically finite functions defined in a neighborhood of $Z$ in $\bar X$, and call the image (germ) of such a function $f$ in $\text{Fin}_Z(\bar X)$ the {\it asymptotic expansion} of $f$ at $Z$. Equivalently, if $\text{Fin}_Z(N_Z\bar X)$ denotes the space of germs, at $Z$, of functions of the form \eqref{equation-asymptotically-finite} defined in a neighborhood of $Z$ in $N_Z \bar X$, the asymptotic expansion of $f$ is the induced map
 $$ \text{Exp}_Z(\bar X) \to \text{Fin}_Z (N_Z X)$$
 from germs of approximate exponential functions (see Definition \ref{definition-approximate-exponential}), which is equivariant for the group bundle $\text{Exp}_Z(N_Z\bar X)$. 
  
  The characters $\lambda$ in an expansion \eqref{equation-asymptotically-finite} will always be assumed to be such that no quotient of two of them extends to a smooth function on $U'$. 
 Under that assumption, the {\it dominant term} of an asymptotically finite function of the form \eqref{equation-asymptotically-finite} is the sum $f_Z=\sum_{\lambda} f_{\lambda} \in C^\infty(U')$; when $U'$ contains the entire orbit of $z$, $f_{\lambda}$ extends uniquely as a generalized $A_Z$-eigenfunction to $X_Z$, and we will consider the dominant term as a function on $X_Z$. (This depends on the orbit of $z$, not just the isomorphism class of $X_Z$!) We write $f\sim f_Z$ to indicate that $f_Z$ is the dominant term of $f$.
\end{definition}

\begin{remark}
 \label{remark-asymptotically-finite-padic}
Notice that, in the non-Archimedean case, the functions $h_\lambda$ in the asymptotic expansion \eqref{equation-asymptotically-finite} are not needed, since they are constant in a neighborhood of $z$; hence, an asymptotically finite function is exactly equal to an $A_Z$-eigenfunction in a neighborhood of $z$.
\end{remark}



\begin{lemma}
 \label{lemma-dominant-term}
The dominant term $f_Z$ of an asymptotically finite function along a $G$-orbit is independent of the choice of an approximate exponential function used to define it.
\end{lemma}

\begin{proof}

[Easy; will be added.]
\end{proof}

In the real case, asymptotically finite functions with respect to a given compactification $\bar X$, set $E$ of ``exponents'' $\lambda$, and bounded degree for the generalized eigenfunctions $f_\lambda$ have a natural structure of a Fr\'echet space. [Details are left to the reader, for now.]



\begin{remark}
\label{remark-expectedtheorem-asymptotics} 
The following is expected to be true for every spherical variety:

\underline{\textbf{Expected theorem:}}

\emph{
Let $X$ denote the points of a homogeneous spherical $G$-variety, and let $X_\Theta$ be a boundary degeneration.\\
  If $\pi$ is any smooth representation of $G$, in the $p$-adic case, and an admissible SF representation of $G$, in the real case, then for any morphism $\ell: \pi\to C^\infty(X)$, there is a unique morphism $\ell_\Theta: \pi\to C^\infty(X_\Theta)$, such that 
  $\ell(v) \sim \ell_\Theta(v)$ for all $v\in \pi$.} (In particular, in the admissible case, $\ell(v)$ is asymptotically finite.)
  
In fact, one can make a stronger statements, where the neighborhood of infinity, or the rate of convergence of asymptotic expansions, is determined by a compact open subset by which $v$ is invariant, resp.\ a continuous seminorm of $v$. This theorem has not appeared in the literature in complete generality. In the next subsections we will formulate (some of) the cases that are known.
\end{remark}

\subsection{Asymptotics in the non-Archimedean case}
\label{subsection-asymptotics-nonarchimedean}



\begin{theorem}
 \label{theorem-asymptotics-nonarchimedean}
Let $X$ denote the points of a homogeneous spherical $G$-variety over a non-Archimedean field, and let $X_\Theta$ be a boundary degeneration. Under the following assumptions:
\begin{itemize}
 \item $G$ is split and $X$ is of \emph{wavefront} type (see \cite[\S 2.1]{Sakellaridis-Venkatesh}), OR
 \item $X$ is symmetric,
\end{itemize}
the following is true: There is a unique morphism
$$ e_\Theta: \mathcal S(X_\Theta)\to \mathcal S(X)$$
with the property that, whenever $X_\Theta$ is realized in the normal bundle of an orbit $Z$ in a smooth toroidal compactification of $X$, and $\phi$ is an approximate exponential map (Definition \ref{definition-approximate-exponential}), for every open compact subgroup $J$ there is a $J$-stable neighborhood $U_Z'$ of $Z$ as in Proposition \ref{proposition-exponential-p-adic}--- in particular, $\phi$ induces a bijection $U_Z'/J=U_X'/J$, where $U_X'$ is the image of $U_Z'$ in $X$--- such that, for $f\in \mathcal S(U_Z')^J$, $e_\Theta(f) = \phi_*(f)$, its pushforward to $U_X'/J$ through this identification.

In particular, the adjoint morphism $e_\Theta^*: C^\infty(X)\to C^\infty(X_\Theta)$ has the property that $e_\Theta^*f|_{U'_Z} = \phi^* f|_{U'_Z}$, for every $f\in C^\infty(X)^J$. 
\end{theorem}

The theorem is expected to hold without these assumptions on $X$.

In particular, if $\ell:\pi\to C^\infty(X)$ is any morphism of smooth representations, we obtain the asymptotic morphism $\ell_\Theta$ of the ``Expected Theorem'' of Remark \ref{remark-expectedtheorem-asymptotics} as $\ell_\Theta = e_\Theta^*\circ \ell$.

\begin{proof}
 See \cite[Theorem 5.1.1]{Sakellaridis-Venkatesh} and \cite[Theorem 1]{Delorme-Plancherel-padic}. 
\end{proof}




\subsection{Asymptotics in the real case}
\label{subsection-asymptotics-real}

\begin{theorem}
 \label{theorem-asymptotics-real}
\begin{enumerate}
 \item Let $X=H$, a (connected) reductive group over $\mathbb R$, under the $G=H\times H$-action. Let $\tau$ be an admissible smooth Fr\'echet representation of moderate growth of $H$, and $\tilde\tau$ its contragredient. Then, for every class $P$ of parabolics in $H$, there exists a finite set $E$ of $A_P$-exponents and a degree $d$, depending on $\tau$, such that all matrix coefficients
 $$ f_{v,\tilde v}(g):=\left< \tau(g) v, \tilde v\right>$$
 are asymptotically finite with exponents $\lambda \in E$ and degree bounded by $d$ in a neighborhood of $P$-infinty, and the map from $\tau\hat\otimes\tilde\tau$ to the corresponding Fr\'echet space $\text{Fin}_P^{E,d}$ of asymptotic expansions is continuous.
 
 In particular, considering only leading terms, there is a morphism $\ell_P:\tau\hat\otimes\tilde\tau \to C^\infty(X_P)$ such that $f_{v,\tilde v} \sim \ell_P(v\otimes\tilde v)$ in a neighborhood of $P$-infinity.
 
 Moreover, if $\ell_P=0$ (i.e., the matrix coefficients of $\tau$ are of rapid decay), for any $P$, then $\tau=0$.
 
 \item Let $X$ be any real spherical variety for a reductive group $G$, and $\pi$ an admissible representation with a \emph{tempered} morphism $\ell: \pi \to C^\infty_{\text{temp}}(X)$, and let $X_\Theta$ denote a boundary degeneration, identified with the open $G$-orbit in the normal bundle of some orbit in a toroidal compactification. Then, there exists a tempered morphism $\ell_\Theta: \pi \to C^\infty_{\text{temp}}(X_\Theta)$, an $A_\Theta$-eigenfunction $h$ on $X_\Theta$ with real positive eigencharacter which is $<1$ on $\exp(\mathfrak a_\Theta^+)$, and a continuous seminorm $q$, such that, for any approximate exponential map $\phi$, $|\phi^* \ell(v) - \ell_\Theta(v)| \le  h\cdot q(v)$ in a neighborhood of $\Theta$-infinity. 
\end{enumerate}

\end{theorem}


\begin{proof}
 For the group case, see \cite[4.4]{Wallach-RR1}. For the tempered case, see \cite{Delorme-Kroetz-Souaifi}.
\end{proof}

\begin{definition}
 \label{definition-asymptotic-matrix-coefficient}
Let $\tau$ be an arbitrary smooth representation of a $p$-adic reductive group $H$, or an admissible smooth representation of moderate growth of a real reductive group $H$. For every class $P$ of parabolics in $H$, let $H_P$ be the corresponding boundary degeneration. The {\it asymptotic matrix coefficient} morphism associated to $P$ is the morphism 
$$m_P: \tau\hat\otimes\tilde\tau \to C^\infty(H_P),$$
where $m_P = \ell_P$ in the notation of Theorem \ref{theorem-asymptotics-real}, in the real case, and $m_P = e_P^*\circ m$, where $m$ is the matrix coefficient map, and $e_P^*: C^\infty(H)\to C^\infty(H_P)$ is the asymptotics map of Theorem \ref{theorem-asymptotics-nonarchimedean}, in the $p$-adic case.
\end{definition}



\section{Consequences of the asymptotics}
\label{section-consequences-asymptotics}

\subsection{Supercuspidals}
\label{subsection-supercuspidals}


\begin{proposition}
 \label{proposition-compact-matrix-coefficients}
For an admissible smooth representation of a real or $p$-adic Lie group $H$, the following are equivalent:
\begin{enumerate}
 \item The matrix coefficients of $\tau$ are of rapid decay (in the real case) or compactly supported (in the $p$-adic case) modulo the center.
 \item The asymptotic matrix coefficient morphisms $m_P$ (Definition \ref{definition-asymptotic-matrix-coefficient} are zero for every class $P$ of proper parabolics in $H$. 
\end{enumerate}

 In particular, in the real case, if the matrix coefficients are of rapid decay modulo the center, then $\tau=0$.
\end{proposition}

\begin{proof}
 Follows immediately from Theorems \ref{theorem-asymptotics-nonarchimedean} and Theorem \ref{theorem-asymptotics-real}, together with the fact that, in the real case, if the asymptotic expansion at infinity is zero, then the function is of rapid decay (modulo center).
\end{proof}



\begin{definition}
 \label{definition-supercuspidal}
Let $H$ be a $p$-adic reductive group. An irreducible admissible representation $\tau$ of $H$ is called {\it supercuspidal} if its matrix coefficients are compactly supported modulo the center.
\end{definition}

\subsection{The subrepresentation theorem}
\label{subsection-subrepresentation}

\begin{theorem}
 \label{theorem-subrepresentation}
Any irreducible admissible representation $\tau$ of a real reductive group $H$, is infinitesimally equivalent to a submodule of an irreducible representation induced from a minimal parabolic; that is, there exists an irreducible (finite-dimensional, necessarily) representation $\sigma$ of the Levi quotient $L$ of the minimal parabolic subgroup $P$ of $H$, and an embedding of $(\mathfrak g, K)$-modules $\tau_{K\text{-fin}} \hookrightarrow I_P(\sigma)_{K\text{-fin}}$, where $I_P(\sigma) = \text{Ind}_P^H(\sigma \delta_P^\frac{1}{2})$ is the (normalized) induced representation.
\end{theorem}

\begin{proof}
 This relies on the statement of Theorem \ref{theorem-asymptotics-real}, that the asymptotics of matrix coefficients in any direction have to be nontrivial. In particular, for the minimal direction we have a non-zero map, which by irreducibility has to be an embedding, $\tau\otimes\tilde \tau \to C^\infty(H_P) = I_{P\times P^-} C^\infty(L)$, whose image consists of $A_P$-finite functions. By projecting to an $A_P$-eigenquotient of the image, we may assume that the image is in an eigenspace, with respect to some character $\chi$ of $A_P$. Notice that $L/A_P$ is compact; hence, the space $C^\infty(L/A_P, \chi)$ has a dense subspace of $L$-finite vectors, which are spanned by matrix coefficients of irreducible representations. Thus, restricting to $K$-finite vectors, there is an morphism (necessarily an embedding) of $(\mathfrak g, K)$-modules $(\tau\otimes\tilde\tau)_{K\times K\text{-fin}} \hookrightarrow I_{P\times P^-} (\sigma\otimes\tilde\sigma)_{K\times K\text{-fin}} = I_P(\sigma)_{K\text{-fin}} \otimes I_{P^-}(\tilde\sigma)_{K\text{-fin}}$, for some irreducible representation $\sigma$ of $L$, and by fixing a vector in $\tilde\tau_{K\text{-fin}}$, we get the embedding claimed in the theorem.
\end{proof}



\section{The Langlands classification}
\label{section-Langlands-classification}

\begin{definition}
 \label{definition-P-dominant}
 Let $G$ be a reductive real or $p$-adic group, let $P\to L$ be a parabolic subgroup with its Levi quotient, and let $\nu:L \to \mathbb C^\times$ be a character. We will say that $\nu$ is {\it $P$-dominant} if $\log(\nu)\in \mathfrak a_P^{*,+}$, and {\it strictly $P$-dominant} if $\log(\nu)\in \mathfrak a_P^{*,+}$. Here, $\mathfrak a_P^* = \Hom(L, \mathbb G_m) \otimes \mathbb R$, $\log$ is the map that sends the absolute value of an algebraic character to its image in $\mathfrak a_P^*$, and $\mathfrak a_P^{*,+}$, $\mathfrak a_P^{*,+}$ are those characters which are non-negative (resp.\ strictly positive) on coroots $\check \alpha$ corresponding to roots in the unipotent radical of $P$, i.e., $|\nu(e^{\check\alpha}(x))| = |x|^\epsilon$ for some $\epsilon \ge 0$ (resp.\ $\epsilon >0$). 
\end{definition}

Equivalently, $\mathfrak a_P^*$ is identified with a subspace of $\mathfrak a^*$ (spanned by the $F$-rational characters of the universal Cartan), and $\mathfrak a_P^{*,+}$ (resp.\ $\mathfrak a_P^{*,+}$) is just the corresponding wall (resp., relative interior of the wall) of the dominant Weyl chamber.


\begin{theorem}[The Langlands quotient theorem]
 \label{theorem-Langlands-quotient}
Let $G$ be a reductive real or $p$-adic group, let $P\to L$ be a parabolic subgroup with its Levi quotient, and let $\tau$ be an irreducible tempered representation of $L$. For any character $\nu:L \to \mathbb C^\times$ which is strictly $P$-dominant, the (normalized) induced representation $I_P^G(\tau\nu)$ has a unique irreducible quotient $\pi_{P,\tau\nu}$, and every irreducible representation of $G$ is of this form, for a unique (up to conjugacy) pair $(P, \tau\nu)$. Moreover, $\pi_{P,\tau\nu}$ is the image of the standard intertwining operator $M_{P^-|P}(\tau\nu): I_P(\tau\nu)\to I_{P^-}(\tau\nu)$. 
\end{theorem}

\begin{proof}

 [Later]
\end{proof}

\begin{example}
 \label{example-trivial-representation}
The trivial representation, for a quasisplit group, is equal to $\pi_{B, \delta^\frac{1}{2}}$, where $B$ is a Borel subgroup, and $\delta$ is its modular character.
\end{example}


\begin{remark}
 \label{remark-Langlands-classification}
The Langlands quotient theorem reduces the classification of irreducible representations to the case of irreducible tempered representations, offering an invaluable link between the ``smooth'' and the ``$L^2$ theory/Plancherel formula'' of irreducible representations. It is also supposed to be compatible with the parametrization provided by the local Langlands conjecture: If $\phi_\tau: \Gamma\to {^LL}$ and $\phi_\nu: \Gamma\to {^LL}$ are Langlands parameters for $\tau$ and $\nu$ (where $\Gamma$, here, denotes the appropriate version of the Weil, or Weil--Deligne group), then $\phi_\tau\cdot\phi_\nu: \Gamma\to {^LL}\hookrightarrow {^LG}$ is a Langlands parameter for $\pi_{P,\tau\nu}$. (Notice that $\phi_\tau$ and $\phi_\nu$ commute, because $\nu$ is a character.)

For example, the Langlands parameter (or rather, its projection to $\check G$) of the trivial representation of a quasisplit group is given by $\Gamma \to \mathbb C^\times \to \check G$, where $\Gamma \to \mathbb C^\times$ is the ``cyclotomic''/absolute value character, and $\mathbb C^\times \to \check G$ is given by $e^{2\rho}: \mathbb G_m\to \check A\subset\check G$ (where $\check A$ is the dual of the universal Cartan).
\end{remark}

\section{Algebraic theory for general totally disconnected groups; the Hecke algebra}
\label{section-td}

In this section, we discuss properties of smooth representations of a general totally disconnected, locally compact group $G$. The category of smooth representations of $G$ will be denoted by $\mathcal M(G)$. We mostly follow the notes from Bernstein's Harvard course \cite{Bernstein-padic}.


\subsection{The Hecke algebra}
\label{subsection-Hecke-algebra}

\begin{definition}
	\label{definition-Hecke-algebra}
The {\it Hecke algebra} $\mathcal H(G)$ of a totally disconnected, locally compact group $G$ is the convolution algebra of compactly supported, complex-valued, smooth measures on $G$ (where ``smooth'' is defined with respect to the left or, equivalently, the right action of $G$ on itself). The Hecke algebra has a natural rational structure $\mathcal H(G)_{\mathbb Q}$, consisting of measures of the form $f dg$, where $f$ is a rational-valued element of $C_c^\infty(G)$, and $G$ is a Haar measure assigning rational volumes to open compact subgroups.

More generally, assume that $G$ is unimodular, and fix a compact open subgroup $J\subset G$. Let $\Lambda$ be a commutative ring. The {\it Hecke algebra of level $J$ with coefficients in $\Lambda$}, $\mathcal H(G,J)_\Lambda$, is the convolution algebra of $\Lambda$-valued, $J$-biinvariant, compactly supported functions on $G$, with convolution defined by assigning Haar measure $1$ on $J$. 
\end{definition}

The Hecke algebra is not unital (unless the group is discrete), but it is idempotented. We recall the what this means. 

\begin{definition}
	\label{definition-idempotented-algebra}
An algebra $A$ is an {\it idempotented algebra} if for every finite collection $\{f_i\}_{i\in I}$ of elements there is an idempotent $e\in A$ with $ef_i = f_ie$ for all $i\in I$. A module $M$ for an idempotented algebra is called {\it nondegenerate} if the natural map 
\begin{equation}
\label{equation-tensor-action}
A\otimes_A M \to M
\end{equation}	
is a surjection (and hence\footnote{The ``hence'' statement uses the property of being idempotented!} an isomorphism); equivalently, if for every $m\in M$ there is an idempotent $e$ with $e\cdot m = m$. 
\end{definition}

We will denote by $\mathcal M(A)$ the category of nondegenerate modules of an idempotented algebra $A$.
The Hecke algebra acts on any smooth $G$-module (see Section \ref{representationtheory-section-measures}), and we have the following.

\begin{proposition}
	\label{proposition-equivalence-Hecke}
	The natural functor gives rise to an equivalence between the categories $\mathcal M(G)$ of smooth $G$-representations and $\mathcal M(\mathcal H(G))$ of nondegenerate $\mathcal H(G)$-modules.
\end{proposition}

\begin{proof}
	When $A=\mathcal H(G)$, the inverse functor is given by the left action of $G$ on $A$ under the isomorphism \eqref{equation-tensor-action}.
\end{proof}

\begin{theorem}
	\label{theorem-smooth-abelian}
 The category $\mathcal M(G)$ of smooth $G$-representations is abelian, and has enough projectives.
\end{theorem}

\begin{proof}
By Proposition \eqref{proposition-equivalence-Hecke}, it is identified with the category of nondegenerate modules for the Hecke algebra. The category of nondegenerate modules for any idempotented algebra is abelian with enough projectives (e.g., we have an epimorphism $\bigoplus_{m\in M} Ae_m\to M$, where $e_m$ is an idempotent fixing $m$. The map sends $e_m$ to $m$, and note that $Ae_m$ is a projective $A$-module, since it represents the exact functor $M\mapsto M^{e_m}=e_m M$.  
\end{proof}

We will show later, using the contragredient (Definition \ref{definition-contragredient}), that -- at least when $G$ is a $p$-adic reductive group -- the category also has enough injectives. For this, we need to know that irreducibles are admissible. 

We have a version of Schur's lemma for irreducible representations of second countable totally disconnected groups.

\begin{proposition}
	\label{proposition-Schur}
If $G$ is a totally disconnected, locally compact, second countable group, and $V \in \mathcal M(G)$ is irreducible, then $\rm{End}_G(V) = \mathbb C$.
\end{proposition}

\begin{proof}
	Since $G$ is second countable, the Hecke algebra has (at most -- we will omit ``at most'' from the rest of the proof) countable dimension, therefore $V=\mathcal H(G)v$ (where $v$ is any nonzero vector) has countable dimension. Therefore, its vector space endomorphisms and therefore its $G$-endomorphisms have countable dimension. Let $T\in \rm{End}_G(V)$. By irreducibility, the operators $(T-\lambda I)$, $\lambda\in \mathbb C$, are either zero or isomorphisms. If one of them is zero, we are done. Otherwise, for dimension reasons, their inverses are linearly dependent, i.e., there is a relation of the form 
	\[ \sum_{i=1}^n c_i(T-\lambda_i I)^{-1}=0 \iff p(T)\prod_{i=1}^n (T-\lambda_i I)^{-1}=0\]
	for some nonzero polynomial $p$. Therefore, $T$ has an eigenvalue $\lambda$, and by irreducibility $T-\lambda I=0$.
\end{proof}

\begin{remark}
 \label{remark-proof-uncountable}
As is obvious from the proof, the analog of the theorem holds when $\mathbb C$ is replaced by any uncountable algebraically closed field $k$ in characteristic zero (for modules of the $k$-rational form of the Hecke algebra over $k$). But we really need the size of the field to play against the size of the group for the conclusion to be true in this generality, for otherwise, an inclusion $k\subset K$ of countable fields would provide a counterexample, with the discrete group $G=K^\times$ acting on the irreducible $k$-vector space $K$, with $G$-endomorphism ring equal to $K$. 
\end{remark}

\subsection{Compact representations}
\label{subsection-compact-representations}

The most important and involved notion that we will encounter in the general theory of totally disconnected groups is that of a compact representation.

\begin{definition}
	\label{definition-compact-representation}
A smooth representation $V\in \mathcal M(G)$ is {\it compact} if its smooth matrix coefficients\footnote{We use $\boxtimes$ to signify that we have a $G\times G$-action on the tensor product; the matrix coefficient map is $G\times G$-equivariant.}
\begin{equation}
	\label{equation-matrix-coefficient}
\tilde V\boxtimes V\to C^\infty(G),
\end{equation}
$\tilde v\otimes v \mapsto (g\mapsto \langle gv, \tilde v\rangle)$, are compactly supported. It is said to be {\it compact modulo center} if the support of (smooth) matrix coefficients is compact modulo the center of $G$.
\end{definition}

Note that, at this point, we do not require the representation to be irreducible, or even finitely generated.

\begin{proposition}
	\label{proposition-compact-admissible}
	An equivalent characterization of compact (resp.\ compact modulo center) representations is that, for any $h\in \mathcal H(G)$ and any $v\in V$, the map $G\ni g\mapsto hgv\in V$ is compactly supported (resp.\ compactly supported modulo center). 
	
	A finitely generated compact representation is admissible.
\end{proposition}

\begin{proof}
	Assume the condition of the proposition, and choose $\tilde v\in\tilde V$, as well as an idempotent $e$ fixing it, then 
	\[\langle gv, \tilde v\rangle = \langle gv, e\tilde v\rangle = \langle e^\vee g v,\tilde v \rangle\]
	 is compactly supported (where by $h\mapsto h^\vee$ we denote the anti-involution $h^\vee(g)=h(g^{-1})$). Vice versa, if there is a pair $(h,v)$ and sequence $g_i \in G$ without compact closure such that $hg_i v\ne 0$, then we can find a linear functional $l\in V^*$ that does not vanish on any of the countably many vectors $hg_i v$. Then the smooth matrix coefficient $\langle g v,  h^\vee l\rangle$ is not compactly supported.
	 
	 If $V$ is generated by a finite set $\{v_i\}_{i\in I}$ of vectors, $J$ is open compact and $e_J$ is the corresponding idempotent (the probability Haar measure on $J$), then $V^J$ is spanned by the vectors $e_J g v_i$ ($g\in G$, $i\in I$), and, by what we just proved, there are finitely distinct such vectors.
\end{proof}


\begin{theorem}
	\label{theorem-compact-splits}
  Assume $G$ to be unimodular (i.e., left Haar measures are also right Haar measures), and let $V$ be a finitely generated, compact representation. If $I_V=\mathcal H(G)_V$ denotes the anti-involution $h\mapsto h^\vee$ applied to the image of the matrix coefficient map \eqref{equation-matrix-coefficient} multiplied by a Haar measure, then $I_V$ is an idempotented two-sided ideal in the Hecke, which admits a complementary idemptotented two-sided ideal $J_V$. This gives rise to a product decomposition $\mathcal M(G)=\mathcal M(I_V)\times \mathcal M(J_V)$ of the category of smooth $G$-representations, i.e., every object $W$ decomposes as $I_VW \oplus J_V W$, and there are no nonzero morphisms between $I_VW$ and $J_V W$.
  
  If $V$ is irreducible, then every object in $\mathcal M(I_V)$ is totally decomposable (into a direct sum of copies of $V$).
\end{theorem}

\begin{proof}
	We will fix a Haar measure in order to consider $\mathcal H(G)$ as a subspace of $L^2(G)$. We claim that orthogonal projection onto the closed span of $I_V$ restricts to a map $\mathcal H(G)\to I_V$. Indeed, this is a $G\times G$-equivariant map, and it is enough to check this claim on $K\times K$-invariants, for any compact open subgroup $K$. But, by Proposition \ref{proposition-compact-admissible}, $V$ is admissible, so $I_V^{K\times K}$ is a finite-dimensional subspace of $L^2(G)^{K\times K}$, hence closed. Therefore, orthogonal projection indeed lands onto $I_V^{K\times K}$. 
	
	Thus, if $J_V$ is the kernel of the orthogonal projection, then $\mathcal H(G)=I_V\oplus J_V$. Both summands are $G\times G$-invariant, hence ideals (since the matrix coefficient map and the orthogonal projection are $G\times G$-equivariant), and if $\{f_i\}{i=1}^n$ is a finite collection of elements in $I_V$, and $e \in \mathcal H(G)$ is an idempotent fixing them under left and right multiplication, with decomposition $e=e_I\oplus e_J$ under the above direct sum, then it is immediate to check that $e_I$ is also an idempotent which fixes them; hence, $I_V$ is idempotented (and similarly for $J_V$). The decomposition of the category is now straightforward.
	
	If $V$ is irreducible, and $W\in \mathcal H(I_V)$ is any module in the $V$-subcategory, then we need to prove that $W$ is generated by the irreducible copies of $V$ inside of it. But, by irreducibility, the matrix coefficient map is an embedding in this case, so, up to the choice of Haar measure, and because of the involution applied, $I_V=V\boxtimes\tilde V$. This is a direct sum of copies of $V$ (as a $G$-module under left multiplication), and therefore $W=I_VW$ is generated by its irreducibles.
\end{proof}

\begin{remark}
	\label{remark-rationality-Hecke}
	At least when $V$ is irreducible, Theorem \ref{theorem-compact-splits} can be proven without recourse to an inner product, hence over other fields $k$ of coefficients. Indeed, the matrix coefficient map (composed with the anti-involution $g\mapsto g^{-1}$) multiplied by a Haar measure and composed with the action of the Hecke algebra gives
	\[V\boxtimes\tilde V\twoheadrightarrow I_V \hookrightarrow \mathcal H(G)\to V\boxtimes\tilde V = (\rm{End}_k V)^\infty,\]
	where the exponent ``infinity'' means smooth under the $G\times G$-action, and the last isomorphism follows from admissibility (Proposition \ref{proposition-compact-admissible}). Now, the composition of these maps cannot be zero (because $V$ is the only irreducible where elements of $I_V$ can act nontrivially), and the codoimain is an irreducible smooth $G\times G$-representation, hence the composition has to be surjective. But then, for every compact open subgroup $K$, the identity of $\rm{End}_k V^K$ is represented by a $K\times K$-fixed element $f_K\in I_V^{K\times K}$, and $J_V^{K\times K}$ can be constructed as the kernel of the action of $f_K$ (on the left or right) on $\mathcal H(G)^{K\times K}$. For more details, see \cite[I.5]{Bernstein-padic}.
\end{remark}

\subsection{Restriction and induction}
\label{subsection-restriction-induction}

We finish this section with generalities regarding induction and restriction.

\begin{definition}
	\label{definition-restriction-induction}
If $H\subset G$ is a closed subgroup of a totally disconnected group, we define the {\it restriction} functor $\rm{Res}^G_H:\mathcal M(G)\to \mathcal M(H)$ in the obvious way, and the (smooth) {\it induction} functor $\rm{Ind}_H^G:\mathcal M(H)\to \mathcal M(G)$ as its right adjoint.
\end{definition}	

Whenever obvious from the context, we will be omitting the exponent $G$ from the notation for induction.

\begin{proposition}
    \label{proposition-induction-exists}
The induction functor is well-defined, and $\rm{Ind}_H^G(V)$ is represented by the space of smooth\footnote{``Smooth'' means in the sense of smooth representations, i.e., any such function is invariant by some compact open subgroup; see Footnote \ref{footnote-smooth}.} functions $f:G\to V$ with the property that $f(hg) = hf(g)$ for every $g\in G$, $h\in H$, with the isomorphism
$$ \Hom_{\mathcal M(G)} (W,\rm{Ind}_H^G(V)) \xrightarrow{\sim}  \Hom_{\mathcal M(H)} (\rm{Res}_H^G(W),V)$$
given by composition with evaluation of such a function $f$ at $1$.
\end{proposition}

\begin{proof}
	The standard proof of Frobenius reciprocity.
\end{proof}

\begin{definition}
	\label{definition-compact-induction}
If $H\subset G$ is a closed subgroup, we define the {\it compact induction} functor $\rm{cInd}_H^G: \mathcal M(H)\to \mathcal M(G)$ as $V\mapsto $ the space of smooth functions $f:G\to V$, with compact support modulo $H$ on the left, and with the property that $f(hg) = hf(g)$ for every $g\in G$, $h\in H$.
\end{definition}

\begin{proposition}
	\label{proposition-induction-restriction-exact}
	All functors $\rm{Res}$, $\rm{Ind}$, $\rm{cInd}$ are exact. If $H\subset G$ is (closed and) open, then $\rm{cInd}_H^G$ is left adjoint to restriction. 
\end{proposition}

\begin{proof}
   Exactness is obvious from the explicit models of these functors. 
   
   If $H\subset G$ is open, then the Hecke algebra of $H$ is a subalgebra of the Hecke algebra of $G$, and one can easily check that $\rm{cInd}_H^G(V)=\mathcal H(G)\otimes_{\mathcal H(H)} V$, so adjunction follows from the equivalence of smooth representations with nondegenerate modules for Hecke algebras (Proposition \ref{proposition-equivalence-Hecke}) and the universal property of tensor products.
\end{proof}

The following two lemmas are best understood in the context of reductive groups, which possess parabolic subgroups $P=LU$ with $P\backslash G$ compact and $U$ (the unipotent radical of $P$) unipotent.

\begin{lemma}
	\label{lemma-compact-induction}
	If $H\subset G$ is a closed subgroup with $H\backslash G$ compact, then the induction functor $\rm{cInd}_H^G = \rm{Ind}_H^G$ preserves admissibility. 
\end{lemma}

\begin{proof}
 For any compact open subgroup $K$, the set of $H\backslash G/K$ double cosets is finite, and for any $K$-invariant $f\in \rm{Ind}_H^G(V)$ and $g\in G$ the evaluation $f(g)$ has to be $H\cap gKg^{-1}$-invariant. Since $H\cap gKg^{-1}$ is open in $H$, the result follows.
\end{proof}
	

\begin{lemma}
	\label{lemma-compact-coinvariants}	
	If $G$ is the union of an increasing sequence of compact open subgroups, $G = \bigcup_{i=1}^\infty G_i$, then the coinvariant functor $\mathcal M(G)\to \rm{Vect}$ sending $V$ to $V_G=$ the quotient of $V$ by the span of vectors of the form $v-gv$, $v\in V$, $g\in G$, is exact.
\end{lemma}

\begin{proof}
	The coinvariant functor is always right exact, for any closed subgroup $H$. If $G$ is compact, it is also left exact, since it can be identified with the invariant functor by integration, $V_G\to V^G$, 
	$$v\mapsto \int_G gv \, dg$$
	(where $g$ is the probability Haar measure on $G$). In particular, in that case, the kernel of the coinvariant map is the kernel of this integral. If $G$ is filtered by compact open subgroups, and $v\in V$ is in the kernel of the coinvariant map, then it will be in the kernel of the coinvariant map $V\to V_{G_i}$ for one of these compact subgroups, which means that its integral over $G_i$ is zero. Applying this to a short exact sequence $0\to A\to B\to C\to 0$, it follows that if $v\in A$ is in the kernel of the coinvariant map for $B$, it is already in the kernel for $A$.
\end{proof}

\section{Algebraic theory for reductive $p$-adic groups; the Bernstein center}
\label{section-padic-reductive}

In this section, $G$ is the set of points over a local, non-Archimedean field $F$, of a connected reductive group $\mathbf G$ defined over $F$. 

\subsection{Parabolic induction and restriction (Jacquet functor)}
\label{subsection-parabolic-induction-restriction}

Recall that a parabolic subgroup $\mathbf P$ of $\mathbf G$ is an algebraic subgroup such that the quotient $\mathbf P\backslash\mathbf G$ is proper. Hence, at the level of $F$-points, $P\backslash G$ is compact. (We also have that $P\backslash G = (\mathbf P\backslash\mathbf G)(F)$.) The quotient of $\mathbf P$ by its unipotent radical is called a \emph{Levi} (subquotient) of $\mathbf G$. 

\begin{definition}
	\label{definition-parabolic-induction-restriction}
	Given a parabolic $P\subset G$ with unipotent radical $U$ and Levi quotient $L$, the {\it parabolic induction} functor $i_P^G:\mathcal M(L)\to \mathcal M(G)$ sends the inflation to $P$ of a representation of $L$ to its induction from $P$ to $G$,
	$$ i_P^G(V) = \rm{Ind}_P^G(\rm{Infl}_L^P(V)).$$
	(Inflation simply means that we consider $V$ as a representation of $P$ via the quotient map to $L$.)
	
	The {\it parabolic restriction} or {\it Jacquet functor} $r^G_P(V)$ sends $V\mapsto V_U$, the $U$-coinvariants, considered as a module for $L$.
\end{definition}


\begin{proposition}
	\label{proposition-Jacquet-exact}
	Parabolic restriction and induction form an adjoint pair of exact functors.
\end{proposition}

\begin{proof}
	Adjunction follows from Proposition \ref{proposition-induction-exists}. Exactness follows from Proposition \ref{proposition-induction-restriction-exact} and, for parabolic restriction, Lemma \ref{lemma-compact-coinvariants}, given the fact that every unipotent subgroup has a filtration by open compact subgroups.
\end{proof}


For the more analytic part of the story it is convenient to define normalized versions of parabolic induction and restriction, where the representation on the Levi is twisted by the square root of the modular character. 

\begin{definition}
	\label{definition-modular-character}
	The {\it modular character} $\delta_P$ of a parabolic subgroup $P$ is the absolute value of its \emph{algebraic modular character}, which is the character $\mathbf L\to \mathbb G_m$ by which the Levi quotient acts on the determinant of the Lie algebra $\mathfrak u$ of its unipotent radical under the \emph{left} adjoint representation. 
\end{definition}

\begin{proposition}
	\label{proposition-modular-character}
	If $d_Lp$ is a left Haar measure on a parabolic subgroup $P$, then $\delta(p)d_Lp$ is a right Haar measure. 
	
	There are identifications (depending on the same choice of positive scalar) of the space $M^\infty(P\backslash G)$ of smooth measures on the flag variety $P\backslash G$ with the induced representation $i_P^G(\delta)$, and of the space $D^\infty(P\backslash G)$ of smooth half-densities on $P\backslash G$ with the induced representation $i_P^G(\delta^\frac{1}{2})$, where $\delta^\frac{1}{2}$ is the positive square root of the modular character.		
\end{proposition}

\begin{proof}
	Measures on $P\backslash G$ are, locally, smooth functions multiplied by the absolute value of a volume form, and there is an identification, up to scalar, of the line bundle of volume forms on $\mathbf P\backslash \mathbf G$ with the algebraic induction of the algebraic modular character. 
\end{proof}

\begin{definition}
	\label{definition-normalized-parabolic-induction-restriction}
	Given a parabolic $P\subset G$ with unipotent radical $U$ and Levi quotient $L$, the {\it normalized parabolic induction} functor $I_P^G:\mathcal M(L)\to \mathcal M(G)$ sends the inflation to $P$ of a representation $V$ of $L$ to the half-density version of its induction from $P$ to $G$, that is, to the space of smooth half-densities $f$ on $P\backslash G$ values in $V$ and with the property that $f(hg) = hf(g)$ for every $g\in G$, $h\in H$. 
	
	The {\it normalized parabolic restriction} or {\it normalized Jacquet functor} $R^G_P(V)$ sends $V\mapsto V_U \otimes D_P^*$, where $D_P = |\det(\mathfrak p\backslash\mathfrak g)^*|^\frac{1}{2}$ (a complex line) denotes the ``fiber'' of the half-density bundle at the point $P1$ of the flag variety $P\backslash G$.
\end{definition}

These definitions are, up to a choice of scalar as in Proposition \ref{proposition-modular-character}, the same as the definitions
$$ I_P^G(V)= i_P^G(V\otimes\delta_P^\frac{1}{2}),$$
$$ R_P^G(V)= r_P^G(V)\otimes\delta_P^{-\frac{1}{2}}$$
more commonly found in the literature. The benefit of the definitions as given above is that they perform the main function of normalized induction indepenently of choices, namely:

\begin{lemma}
	\label{lemma-normalized-induction-unitary}
	Normalized restriction and induction form an adjoint pair of exact functors, and if $V\in \mathcal M(L)$ is unitary, i.e., endowed with an $L$-invariant Hilbert space (semi)norm, then $I_P^G(V)$ is unitary, with (semi)norm
	$$\Vert f\Vert =  \left(\int_{P\backslash G} \Vert f(g)\Vert \right)^\frac{1}{2}.$$
\end{lemma}

Note that the integral defining $\Vert f\Vert$ makes sense precisely because $f$ is a half-density (valued in $V$).

\begin{proof}
	The first statement follows from Proposition \ref{proposition-Jacquet-exact}, and the fact that the norm is well-defined and $G$-invariant is immediate.
\end{proof}


\subsection{Compact and cuspidal representations}
\label{subsection-compact-cuspidal}

Now we will characterize compact-mod-center representations (Definition \ref{definition-compact-representation}), for reductive $p$-adic groups, in terms of their Jacquet functors.


\begin{definition}
	\label{definition-cuspidal}
A smooth representation $V\in \mathcal M(G)$ is called {\it quasicuspidal} if its Jacquet functors $R_P V$ are zero for all proper parabolic subgroups $P\subset G$, {\it cuspidal} if it is quasicuspidal and finitely generated, and {\it supercuspidal} if it is (quasi)cuspidal and of finite length.
\end{definition}

Please note that there are slightly different preferences in the literature as to how to distinguish among the three terms.

Our main result will be that a smooth representation $V\in \mathcal M(G)$ is quasicuspidal if and only if it is compact modulo center (Theorem \ref{theorem-compact-cuspidal}), but we need some preparation. First, we need the following two results about $p$-adic reductive groups [whose proof should be added sometime!]. Let $P_0\subset G$ be a minimal parabolic subgroup (i.e., the $F$-points of a minimal $F$-rational parabolic), with $A=$ the maximal split torus in the center of a Levi subgroup $L_0$ of $P_0$. (In particular, $L_0$ is the centralizer of $A$, and the choice of $L_0$ determines an opposite parabolic $P_0^-$ with $P_0^-\cap P_0=L_0$.) We let $\Lambda$ denote the cocharacter (weight) lattice of $A$, $\Lambda^+\subset \Lambda$ the submonoid of $P_0$-dominant weights, and will use the notation $a^\lambda$ to denote the image of an element $a\in F^\times$ under $\lambda\in \Lambda$. For $\lambda\in \Lambda^+$, we denote by $L_\lambda$ the centralizer of $\lambda$ (= a Levi subgroup containing $L_0$), and by $P_\lambda$, $U_\lambda$ (resp.\ $P_\lambda^-$, $U_\lambda^-$) the parabolic and unipotent radical with Levi $L_\lambda$ containing $P_0$ (resp.\ $P_0^-$).

\begin{description}
	\item[Weak Cartan decomposition] There is a compact subset $K\subset G$ such that $G = K \varpi^{\Lambda^+} K$, where $\varpi$ is any uniformizer in $F$.
	\item[Basis of Iwahori-type subgroups] There is a sequence of open compact subgroups $J$ which form a basis of neighborhoods of the identity, and possess Iwahori factorization with respect to any standard parabolic $P_\lambda$, that is, $J$ is the product of its subgroups $J_\lambda^0:=J\cap L_\lambda$, $J_\lambda^+=J\cap U_\lambda^+$, and $J_\lambda^-=J\cap U_\lambda^-$, in any order. 
\end{description}



\begin{proposition}
	\label{proposition-Hecke-subalgebra}
Given an subgroup $J$ with Iwahori factorization, there is an embedding of the monoid algebra $\mathbb C[\Lambda^+]$ into the Hecke algebra $\mathcal H(G,J)$, 
\begin{equation}
	\label{equation-Hecke-subalgebra}
	\mathbb C[\Lambda^+]\hookrightarrow \mathcal H(G,J),
\end{equation}
sending the basis element indexed by $\lambda\in \Lambda^+$ to a scalar multiple of the characteristic measure of the double coset $J\varpi^\lambda J$.
\end{proposition}

The precise normalization of this embedding, of course, depends on how one chooses the scalar multiple on generators of the monoid.

\begin{proof}
	This follows immediately from the following calculation of products of double cosets, when $\lambda_1,\lambda_2\in \Lambda^+$.
	$$ J\varpi^{\lambda_1}J \cdot J\varpi^{\lambda_2}J = J\varpi^{\lambda_1}J_0^+ J_0^0  J_0^- \varpi^{\lambda_2}J = J\varpi^{\lambda_1+\lambda_2}J,$$
	since for dominant $\lambda$ we have $\varpi^\lambda J_0^+ \varpi^{-\lambda} \subset J_0^+$ and $\varpi^{-\lambda} J_0^- \varpi^{\lambda} \subset J_0^-$ (and $J_0^0 $ commutes with $\varpi^\lambda$).
\end{proof}

\begin{remark}
	\label{remark-Up-operator}
   For $G=\rm{GL}_2$ over $\mathbb Q$, $J$ the usual Iwahori subgroup $\begin{pmatrix} \mathbb Z_p^\times & \mathbb Z_p \\ p\mathbb Z_p & \mathbb Z_p^\times\end{pmatrix}$ and $\varpi^\lambda = p^\lambda = \begin{pmatrix} p \\ & 1 \end{pmatrix}$, the corresponding element of the Hecke algebra is the well-known {\it $U_p$ operator}.
\end{remark}


\begin{theorem}
	\label{theorem-compact-cuspidal}
 A smooth representation $V\in \mathcal M(G)$ is quasicuspidal if and only if it is compact modulo center. 
\end{theorem}

\begin{proof}
By Proposition \ref{proposition-compact-admissible} and the weak Cartan decomposition, it is enough to show that for any $v\in V$, any open compact subgroup $J$, and $k_1,k_2$ in a fixed compact subset $K\subset G$, the function $\lambda\mapsto J k_1 \varpi^\lambda k_2 v$ is compactly supported modulo center on $\Lambda^+$. (Here and later, by the action of a compact group on a vector we mean the averaging operator, i.e., the action of its probability Haar measure.)


Obviously, $k_2$ can be ignored since it just changes $v$. Moreover, the subset $JK$ is compact open and hence a finite union of right cosets of some (other) compact subgroup $J'$, so, by replacing $J$ with $J'$ we are reduced to showing that the function $\lambda\mapsto J\varpi^\lambda v$ is compactly supported modulo center. Finally, we could have chosen $J$ small enough that $v$ is $J$-invariant, and therefore $J\varpi^\lambda  v = J\varpi^\lambda J v$, i.e., we are considering the action of the subalgebra $\mathbb C[\Lambda^+]$ of Proposition \ref{proposition-Hecke-subalgebra} on $V^J$.

The basic observation, in that case, is the following: The kernel of the Hecke operator $J\varpi^\lambda J$ acting on $V^J$ is equal to the kernel of $\varpi^{-\lambda} J_\lambda^+ \varpi^\lambda$ (on $V^J$). Indeed, this follows from the calculation 
$$J\varpi^\lambda J = \varpi^\lambda \left(\varpi^{-\lambda} J_\lambda^+ \varpi^\lambda\right) J,$$ 
and the fact that $\varpi^\lambda$ acts invertibly. 

As $\Lambda$ varies in $\Lambda^+$, the parabolics $P_\lambda$ vary over the finite set of ``standard'' parabolics containing $P_0$; number them $P_0, \dots, P_N$, with $P_N=G$. If $v\in V^J$ is in the kernel the Jacquet module (coinvariant) map with respect to all $P_i$'s other than $G$, then for each $i<N$ there is a compact open subgroup $U_i^0$ of the corresponding unipotent radical $U_i$ which annihilates $v$ -- fix such a subgroup for every $i$. Then, for $\lambda\in \Lambda^+$ outside of a compact (=finite) subset modulo  the center, $\varpi^\lambda J_\lambda^+ \varpi^{-\lambda}$ contains one of the $U_i^0$'s, which means that $v$ is in the kernel of the Hecke operator $J \varpi^\lambda J$. Thus, if $V$ is quasicuspidal, it is compactly supported modulo the center.

Vice versa, if $v \in V^J$ being in the kernel of $J \varpi^\lambda J$ and hence of $\varpi^{-\lambda} J_\lambda^+ \varpi^\lambda$ implies that $v$ is in the kernel of the Jacquet module for $P_\lambda$. If this happens for all large enough non-central $\lambda$, it is in the kernel of the Jacquet modules for all proper parabolics. Thus, if $V$ is compact modulo center, it is quasicuspidal.
\end{proof}


\subsection{Cuspidal components}
\label{subsection-cuspidal-components}

We will now see that the equivalence class of an irreducible supercuspidal representation modulo appropriate twists splits a very simple component off the category $\mathcal M(G)$ of smooth representations of $G$. The ``twists'' are characters of the discrete abelian group $G/G_0$, where $G_0$ is the subgroup of $G$ generated by all compact subgroups. 

We will give a more algebraic definition of $G_0$. If $\mathbf A^G$ is the maximal split torus quotient of $\mathbf G$, we let $G_0$ be the kernel of the map $G\to \mathbf A^G(F)/\mathbf A^G(\mathfrak o)$. The quotient $G/G_0$ is a (full rank) sublattice of $\mathbf A^G(F)/\mathbf A^G(\mathfrak o)$, although not necessarily equal to it, and we let $\Psi_G$ be the complex algebraic group whose character group is $G/G_0$. We will also abuse notation and write $\Psi_G$ for its group of complex points, which can be identified with the characters $G/G_0\to \mathbb C^\times$. 

\begin{remark}
	\label{remark-unramified-characters}
Bernstein's notes \cite{Bernstein-padic} call $\Psi_G$ the set of unramified characters, but there can be more unramified characters in the number-theoretic sense. For example, if $\chi:F^\times/\mathfrak o^\times\to \{\pm 1\}$ is the quadratic unramified character of $F^\times$, then $\chi\circ\det$ deserves to be called an unramified character of the group $\rm{PGL}_2(F)=\rm{GL}_2(F)/F^\times$, despite the fact that $G=G_0$ for $\rm{PGL}_2$. 
\end{remark}

Consider the set $\rm{Irr}(G)$ of isomorphism classes of irreducible smooth representations of $G$; a priori, this is just the set, but we will see that it admits a finite-to-one map to a certain complex algebraic variety. The group $\Psi_G$ acts on $\rm{Irr}(G)$.

\begin{lemma}
	\label{lemma-finite-stabilizer}
The stabilizer of every $\pi\in \rm{Irr}(G)$ under the $\Psi_G$-action is finite.	
\end{lemma}

\begin{proof}
	If $\mathbf A_G$ denotes the maximal split torus in the center of $\mathbf G$, the map $\mathbf A_G\to \mathbf A^G$ has finite cokernel (and kernel). For $\eta\in \Psi_G$, a necessary condition for $\pi$ and $\pi\otimes\eta$ to be isomorphic is that they have the same central character, which means that $\eta$ can only belong to the finite group of characters of $G/G_0$ which are trivial in the image of $A_G$. 
\end{proof}

\begin{lemma}
	\label{lemma-restriction-G0}
	If $\pi$ is an irreducible representation of $\pi$ then its restriction to $G_0$ is semisimple of finite length. The restrictions of two such representations $\pi$ and $\pi'$ share an irreducible summand iff they belong to the same $\Psi_G$-orbit (in which case the restrictions coincide). 
\end{lemma}

\begin{proof}
	By Schur's lemma (Proposition \ref{proposition-Schur}), the center and, in particular, its maximal split torus $A_G$ acts by a character on $\pi$. The product $H=A_GG_0$ has finite index in $G$, and the induction of the restriction of $\pi$ to $H$ is the direct sum of copies of $\pi$ twisted by characters of $G/H$. From this, it follows that the restriction is semisimple of finite length. (Finite length, because the length can only go up by induction; semisimple, because every $H$-equivariant quotient of $\rm{Res}_H^G(\pi)$ induces a $G$-equivariant quotient of $\rm{Ind}_H^G \rm{Res}_H^G(\pi)$, and the splitting of the latter implies, by projection to the identity component of $G/H$, a splitting of the former.)
	
	Regarding the case where $\pi|_{G_0}$ and $\pi'|_{G_0}$ share an irreducible summand, we can view the space $\Hom_{G_0}(\pi,\pi')$ as a representation of the abelian group $G/G_0$. By finite length, it is finite-dimensional, and therefore has an eigenvector with some eigencharacter $\eta$. But this means that $\pi' = \eta\pi$.
\end{proof}


\begin{definition}
	\label{definition-cuspidal-component}
A {\it cuspidal component} of $\rm{Irr}(G)$ is the $\Psi_G$-orbit of a supercuspidal representation.
\end{definition}

Note that we do not, a priori, know if any cuspidal representations exist at all. But this is not a question that we will address here.



\begin{theorem}
	\label{theorem-cuspidal-component}
Every cuspidal component $D$ splits the category $\mathcal M(G)$ of smooth representations as 
$$\mathcal M(G)=\mathcal M(G)_D \times \mathcal M(G)_{nD},$$
where all irreducible subquotients of $\mathcal M(G)_D$ and no irreducible subquotients of $\mathcal M_{nD}$ belong to $D$.

More generally, we have a decomposition 
\begin{equation}
	\label{equation-cuspidal-decomposition}
	\mathcal M(G)=\prod_D\mathcal M(G)_D \times \mathcal M(G)_{nc},
\end{equation}
where $D$ runs over the set of all cuspidal components, and $\mathcal M(G)_{nc}$ has no supercuspidal subquotients. The product over $D$ corresponds to the subcategory of quasicuspidal representations, and the equivalence (from right to left) is given by the functor of direct sum.
\end{theorem}


\begin{proof}
The restriction of (any) $\pi\in D$ to $G_0$ is a finite direct sum of irreducibles by Lemma \ref{lemma-restriction-G0}, which are cuspidal, since Jacquet modules do not see the difference between $G$ and $G_0$. Therefore, they are compact by Theorem \ref{theorem-compact-cuspidal} (and the fact that $G_0$ has compact center), and each of them splits the category of $G_0$-modules by Theorem \ref{theorem-compact-splits}. Putting the components corresponding to all summands together, the first claim follows.

The second claim will also follow if we prove that, for every compact open subgroup $K$ of $G$, there is only a finite number of cuspidal components $D$ containing $V\in \mathcal M(D)$ with $V^K\ne 0$; indeed, this would allow us to perform the decomposition on the submodules of $V\in \mathcal M(G)$ generated by $V^K$, for all $K$. This follows from the \emph{uniform admissibility theorem} [to be added].
\end{proof}

We will now describe the category $\mathcal M(G)_D$ explicitly as a module category for a finite algebra over the coordinate ring $\mathbb C[D]$. For this, we need some general algebraic preliminaries. Let $\mathcal C$ be an abelian category with arbitrary coproducts. A {\it progenerator} is a \emph{compact, projetive generator} of the category, that is, an object $P\in \mathcal C$ such that the functor 
$$ N\mapsto \Hom(P,N)$$
is \emph{exact, faithful, and commutes with arbitrary coproducts}. Under the projective assumption, ``faithful'' simply translates to the condition that $\Hom(P,N)\ne 0$ for all nonzero objects $N\in \mathcal C$. This functor lands in abelian groups, but since we can precompose any homomorphism with an endomorphism of $P$, it can be upgraded to a functor
\begin{equation}
	\label{equation-Morita-functor}
\mathcal C\to \rm{End}(P)^\circ\text{-modules}.
\end{equation}

\begin{theorem}
	\label{theorem-Morita}
Under the assumptions above, the functor \eqref{equation-Morita-functor} is an equivalence of abelian categories.
\end{theorem}

\begin{proof}
	Omitted. This is often referred to as the {\it Gabriel--Mitchell theorem}. See \cite[Theorem II.1.3]{Bass-Ktheory}.
\end{proof}

Note that there can be many projective generators in such a category; their endomorphism rings are then {\it Morita equivalent}, i.e., their module categories are equivalent. 

One example of a projective generator for $\mathcal M(G)_D$ is 
\begin{equation}
	\label{equation-progenerator-cuspidal}	
	\Pi_D= \rm{cInd}_{G_0}^G \rm{Res}_{G_0}^G \pi,
\end{equation}
where $\pi$ is any element of $D$. (Obviously, all have equal restrictions.) Equivalently,
\[\Pi_D = \pi\otimes \mathbb C[G/G_0],\]
and recall that the group ring $\mathbb C[G/G_0]$ is the coordinate ring of $\Psi_G$. 

Fixing a base point $\pi\in D$, we can identify $D$ with $\Psi_G/\Gamma$, where $\Gamma$ is the (finite, by Lemma \ref{lemma-finite-stabilizer}) stabilizer of the isomorphism class of $\pi$ in $\Psi_G$. We will now define a $2$-cocycle $\sharp:\Gamma \times \Gamma\to \mathbb C^\times$, as follows. For each $\gamma\in\Gamma$, choose an isomorphism $\alpha_\gamma: \pi\xrightarrow\sim \pi\otimes\gamma$. By tensoring with any $\chi\in \Psi_G$ we can also consider $\alpha_\gamma$ as am isomorphism $\pi\otimes\chi\to \pi\otimes\chi\gamma$, and by Schur's lemma these isomorphisms need to satisfy a relation of the form $\alpha_{\gamma_1}\circ \alpha_{\gamma_2} = \sharp(\gamma_1,\gamma_2) \alpha_{\gamma_1\gamma_2}$ for all $\gamma_1, \gamma_2 \in \Gamma$. This is our $2$-cocycle, whose cohomology class does not depend on the isomorphisms chosen.

Given such a $2$-cocycle, we can define a {\it twisted group algebra} $\mathbb C[\Gamma, \sharp]$ with basis $T_\gamma$, $\gamma\in \Gamma$ and relations $T_{\gamma_1}T_{\gamma_2} = \sharp(\gamma_1,\gamma_2) T_{\gamma_1\gamma_2}$. The {\it crossed product}
\begin{equation}
	\label{equation-crossed-product}
\mathbb C[\Psi_G]\ltimes \mathbb C[\Gamma,\sharp]
\end{equation}
is the tensor product of the two factors, as a vector space, with multiplication
$$ (f_1\otimes T_{\gamma_1})(f_2\otimes T_{\gamma_2}) = (f_1 \cdot {^{\gamma_1}f_2})\otimes \sharp(\gamma_1,\gamma_2) T_{\gamma_1\gamma_2}.$$
Here, $f^\gamma$ denotes the action of $\gamma$ on $f$ by translation. 

\begin{proposition}
	\label{proposition-endomorphism-cuspidal}
The endomorphism algebra of the progenerator \eqref{equation-progenerator-cuspidal} is the crossed product \eqref{equation-crossed-product}. Its center is the coordinate ring $\mathbb C[D] = \mathbb C[\Psi_G]^\Gamma$.
\end{proposition}

\begin{proof}
	Left as exercise.
\end{proof}

\begin{remark}
	The crossed product algebra is an Azumaya algebra, and is Morita equivalent to its center $\mathbb C[D]$ iff the cohomology class of the cocycle $\sharp$ is trivial; see \cite[Theorem A.15]{Auslander-Goldman}.
\end{remark}


\subsection{Cuspidal support}

We will now extend the decomposition of the category of smooth representations given in Theorem \ref{theorem-cuspidal-component} to a decomposition indexed by cuspidal components of Levi subgroups.

\begin{definition}
	\label{definition-cuspidal-datum}
A {\it cuspidal datum} is a pair $(L,\sigma)$ consisting of a Levi subgroup $L$ of $G$ and an irreducible supercuspidal representation $\sigma$ of $L$. 

Two cuspidal data $(L,\sigma)$ and $(M,\tau)$ are said to be {\it associate cuspidal data} if there is a $g\in G$ such that $gMg^{-1}=L$ and $\sigma(g\bullet g^{-1})$ is isomorphic to $\tau$. Two cuspidal components for Levi subgroups, $(L,D)$ and $(M,D')$ are said to be {\it associate cuspidal components} if one (equivalently, all) elements of the former is associate to an element of the latter.

A {\it Bernstein component} is the following complex variety whose complex points parametrize equivalence classes of cuspidal data: Given a Levi subgroup $L$ and a cuspidal component (Definition \ref{definition-cuspidal-component}) $D$ of $L$ (considered as a complex affine variety), its Bernstein component $\Omega(L,D)$ is the quotient variety $\Omega\sslash W(L,D) = \rm{Spec} \mathbb C[D]^{W(L,D)}$, where $W(L,D)$ is the subgroup of the (finite) group $\mathcal N_G(L)/L$ which maps the component $D$ to itself (i.e., such that precomposition with a representative in $\mathcal N_G(L)$ preserves the component of a cuspidal representation $\sigma\in D$). This is identified with $\Omega(M,D')$ for any associate cuspidal component $(M,D')$ in the natural way (i.e., via conjugation by any element $g\in G$ giving rise to the associate relation). The disjoint union of all Bernstein components is the {\it Bernstein variety} $\Omega_G$.
\end{definition}

There is a countable number of Bernstein components, and therefore the Bernstein variety is a countable union of affine irreducible varieties. This means that its coordinate ring $\mathbb C[G]$ is the \emph{product} of the coordinate rings of its irreducible components. 


We will now start studying parabolically induced cuspidal blocks. It will sometimes be convenient to work with the normalized induction and restriction functors $I_P=I_P^G$, $R_P=R_P^G$ of Definition \ref{definition-normalized-parabolic-induction-restriction}, rather than the unnormalized ones $i_P$, $r_P$ (although for many statements the difference plays no role). Our basic tool will be the {\it Geometric Lemma}, which is the proposition that follows.

\begin{proposition}[Geometric Lemma]
	\label{proposition-geometric-lemma}
Let $P$ and $Q$ be parabolics in $G$ with Levi quotients $L$, $M$, respectively. The functor
$$\Gamma: \mathcal M(L) \to \mathcal M(M)$$
given by $\Gamma = R_Q \circ I_P$ has a finite filtration of by subfunctors indexed by the poset of double cosets $w\in P\backslash G/Q$ (with the opposite to the closure relation order), with quotients naturally isomorphic (up to a scalar) to
$$ \Gamma_w = I^M_{M\cap w^{-1}Pw} \circ w \circ R^L_{L\cap wQw^{-1}}.$$
Here, by abuse of notation, we identify a double coset with a representative in $G$, and the operator ``$w$'' appearing in the formula is simply conjugation by $w$, which turns modules for the Levi quotient of $L\cap wQw^{-1}$ into modules for the Levi quotient of $M\cap w^{-1}Pw$.
\end{proposition}

Note that, for any two parabolics $P$ and $Q$, the interesection $P\cap Q$ projects to parabolic subgroups modulo the unipotent radical of each. An easy way to see this is to remember that the intersection contains a maximal split torus of $G$, and that these parabolics are then of the form $P_\lambda$, for some cocharacters $\lambda$ into that torus, in the notation used at the beginning of Subsection \ref{subsection-compact-cuspidal}. Also, to write intersections such as $M\cap w^{-1}Pw$, we have chosen a spliting of the Levi quotient of $Q$ into $Q$, but this intersection does not depend on the choice of splitting.

\begin{proof}
	Using the explicit description of induction $I_P^G V$ in terms of smooth half-densities on $P\backslash G$ valued in $V$, every $Q$-invariant open subset of $P\backslash G$ gives rise to a filtered piece $F^w I_P^GV$, of half-densities (compactly) supported in that subset. Every $w\in P\backslash G/Q$ defines such an open subset $S_w$, by taking the union of double cosets which contain it in their closure. Moreover -- and this is something particular about totally disconnected spaces -- every such section that vanishes on the $w$-orbit is supported in its (open) complement in $S_w$. Therefore, the quotient $F^w V/\sum_{w'<w} F^{w'} V$ is identified with sections over the $w$-orbit of the sheaf of half-densities on $P\backslash G$, valued in $V$. \emph{For notational simplicity, now, we assume that $w$ is represented by the identity in $G$, as we may by replacing $P$ by $w^{-1}Pw$.}	

	If we choose an identification of the sheaf of half-densities with the induction of $\delta_P^{\frac{1}{2}}$, as a $Q$-module this graded piece is isomorphic to 
	$$ \rm{cInd}_{P\cap Q}^Q \delta_P^{\frac{1}{2}} V.$$
	We can write the unipotent radical $U_Q$ of $Q$ as a product of subgroups $U_Q = (U_Q\cap P) \times  U_Q'$, with the first factor being the stabilizer of the point $P1$ in $P\backslash G$, and the second acting freely on its $Q$-orbit. (Indeed, think of roots with respect to a maximal split torus, as in the comment preceding this proof.) The (unnormalized) $Q$-Jacquet module of this graded piece will then be the composition of averaging over $U_Q'$-orbits, followed by coinvariants with respect to $U_Q\cap P$. The latter, after projection modulo $U_P$, is precisely the unipotent radical of the parabolic $L\cap Q$ of $L$. Similarly, $P\cap Q\backslash Q/U_Q$ is the parabolic $M\cap P$ of $M$. Therefore, for the normalized Jacquet module (again choosing an identification of the sheaf of half-densities on $Q\backslash G$ with the induction of $\delta_Q^{\frac{1}{2}}$), we get a graded piece 
	$$ \Gamma_1 V =  \delta_Q^{-\frac{1}{2}}  i^M_{M\cap P}  r^{L}_{L\cap Q} \delta_{U_Q'} \delta_P^{\frac{1}{2}} V,$$
	where the factor $\delta_{U_Q'} :=$ the modular character for the action of $L\cap Q$ on $U_Q'$ appears because of the $U_Q'$-integration. One now checks that $\delta_P \delta_{U_Q'} = \delta_{M\cap P}$ and $\delta_Q \delta_{U_Q'}^{-1} = \delta_{L\cap Q}$, hence this becomes 
	$$  i^M_{M\cap P} \delta_{M\cap P}^{\frac{1}{2}}  \delta_{L\cap Q}^{-\frac{1}{2}}  r^{L}_{L\cap Q}   V = I^M_{M\cap P} R^L_{L\cap Q} V.$$
\end{proof}

A corollary of the geometric lemma is the following.

\begin{proposition}
	\label{proposition-irreducible-induction}
	Let $\sigma$ be an irreducible, admissible, unitarizable representation of a Levi subgroup $L$. There is an open dense $U\subset \Psi_L$ such that, for $\eta\in U$ and $P$, $Q$ two parabolic subgroups with Levi $L$, the representations $I_P(\sigma\otimes\eta)$ and $I_Q(\sigma\otimes\eta)$ are irreducible and isomorphic.
\end{proposition}

\begin{proof}
Irreducibility is an open condition: Parabolic induction preserves admissibility, and irreducibility of an admissible representation $V$ is equivalent to the statement that the action map $\mathcal H(G,J)\to \rm{End}(V^J)$ is surjective for every $J$. We can choose a compact open subgroup $K$ such that $G=PK$, and then the restrictions of all $V=I_P(\sigma\otimes\eta)$ to $K$ are identified (with the induction of the restriction to $P\cap K$ of $\sigma$). Hence, for all $J\subset K$, the vector spaces $\rm{End}(V^J)$ are identified, and the action of $\mathcal H(G,J)$ gives a map 
$$ \mathcal H(G,J) \to \mathbb C[\Psi_L]\otimes \rm{End}(V^J).$$
Therefore, this map being surjective is an open condition. [For the conclusion, it is enough to restrict to a $J$ of Iwahori type such that these modules are generated by their $J$-invariants -- to be added.]

Therefore, it is enough to prove irreducibility for a single element in the family. If $\eta$ is a unitary character, then $I_P(\sigma\otimes\eta)$ is unitarizable and admissible, which by a standard argument (projection to the orthogonal complement) implies that it is completely reducible. Therefore, it is enough to show that, for some unitary $\eta$, the $G$-endomorphism ring of $I_P(\sigma\otimes\eta)$ is scalar. 

But $\Hom_G(I_P(\sigma\otimes\eta), I_P(\sigma\otimes\eta))= \Hom_L(R_P I_P(\sigma\otimes\eta), \sigma\otimes\eta)$, and, by Proposition \ref{proposition-geometric-lemma}, the functor $R_P I_P$ has a filtration with graded quotients $I^L_{L\cap w^{-1}Pw} \circ w \circ R^L_{L\cap wLw^{-1}}$, with $w$ ranging over $P\backslash G/P$. The point now is that, for $w\ne 1$ and $\eta$ generic in $\Psi_L$ (and hence in its unitary subset), the Jordan--H\"older components of these graded pieces have different central characters than $\sigma\otimes\eta$, and therefore this endomorphism ring is scalar. This proves the first claim. 

[Proof of second claim to be added.]
 
\end{proof}



\begin{proposition}
	\label{proposition-cuspidal-data-equivalences}
Given a Bernstein component $\Omega$ represented by a cuspidal component $D$ of a Levi subgroup $L$, the following are equivalent for a smooth representation $V\in \mathcal M(G)$.
\begin{enumerate}
	\item \label{item:one} $V$ injects into a direct sum, indexed by parabolics $P$ with Levi subgroup $L$, of representations $i_P(W_P)$, with $W_P\in \mathcal M(L)_D$.
	\item $V$ is a subquotient of representations as in \eqref{item:one}.
	\item If $(M,D')$ is any cuspidal datum that is \emph{not} associate to $(L,D)$, and $Q$ a parabolic with Levi $M$, the Jacquet module $r_Q(V)$ has zero component in $\mathcal M(M)_{D'}$.
	\item For a parabolic $P$ with Levi $L$, the Jacquet module $r_P(V)$ belongs to the product of $\mathcal M(L)_{D'}$ with $(L,D')$ associate to $(L,D)$, and if we denote by the index $D$ the projetion to the $D$-cuspidal block, the map $V\mapsto \bigoplus_P i_P(r_P V)_D$, with $P$ ranging over all parabolics with Levi $L$, is injective.
\end{enumerate}
\end{proposition}

\begin{proof}
This follows relatively easily from the Geometric Lemma \ref{proposition-geometric-lemma}. See \cite[Proposition-definition 2.8]{Bernstein-centre} for details.
\end{proof}





\begin{definition}
	\label{definition-Bernstein-block}
Given a Bernstein component $\Omega\subset \Omega_G$, in the sense of Definition \ref{definition-cuspidal-datum}, the corresponding {\it Bernstein block} of the category $\mathcal M(G)$ is the full subcategory $\mathcal M(G)_\Omega$ of all representations $V\in \mathcal M(G)$ satisfying the equivalent conditions of Proposition \ref{proposition-cuspidal-data-equivalences}. 
\end{definition}



\begin{theorem}
	\label{theorem-Bernstein-decomposition}
The category $\mathcal M(G)$ is the product of its Bernstein blocks $\mathcal M(G)_\Omega$, with the equivalence given by direct sums.
\end{theorem}

That means that every object decomposes into a direct sum of objects belonging to the various blocks, and morphisms between two objects are products of the morphisms between their components. 

\begin{proof}
 For every representation $V$ adjunction gives rise to a map 
 $$ V\to \prod_{(L,P,D)} I_P (R_PV)_D,$$
 where $(L,D)$ ranges over a set of representatives for all Bernstein components, and $P$ ranges over all parabolics with Levi $L$; the index $D$ means projection to the corresponding cuspidal component for the Levi. 
 
 The kernel of this map is a submodule with the property that the cuspidal projections of all its Jacquet modules (including for $P=G$) vanish. Such a representation is clearly zero. 
 
 The image lies in the direct sum, by Theorem \ref{theorem-cuspidal-component}.
\end{proof}

\subsection{The Bernstein center}
   \label{subsection-Bernstein-center}

\begin{definition}
	\label{definition-Bernstein-center}
The {\it center of an abelian category} $\mathcal C$ is the ring of endomorphisms of its identity functor. The {\it Bernstein center} $\mathfrak Z(G)$ of a reductive $p$-adic group $G$ is the center of the category $\mathcal M(G)$ of its smooth representations.
\end{definition}

Explicitly, an element $z$ of the Bernstein center consists of a collection $z_V$ of ($G$-)endomorphisms of its various smooth representations $V$, which is functorial, in the sense that for every morphism $V\to W$ the following diagram commutes.
$$ \xymatrix{V \ar[r] \ar[d]_{z_V} &W \ar[d]^{z_W} \\ V \ar[r]& W}.$$
We will ignore general set-theoretic questions, which are pretty easy to settle in module categories; for example:

\begin{example}
	\label{example-center-modulecategory}
	If $A$ is a unital algebra, and $\mathcal C = \rm{Mod}(A)$, the category of $A$-modules, its center is the center of $A$. Indeed, the category is generated by the free object $A$, and any $A$-endomorphism of $A$ is determined by the image of its identity element, which has to be central. This clearly induces functorial $A$-endomorphisms of all modules $N$, through presentations of the form $\bigoplus_I A\to N$.
\end{example}

In our case, $\mathcal M(G)$ is equivalent to the category $\mathcal M(\mathcal H(G))$ of nondegenerate modules of its Hecke algebra (Proposition \ref{proposition-equivalence-Hecke}), which however is only idempotented, not unital. We can build a unital ``completed'' Hecke algebra out of it.

\begin{definition}
   \label{definition-completed-Hecke}
The {\it completed Hecke algebra} $\widehat{\mathcal H(G)}$ is the convolution algebra of {\it essentially compactly supported} distributions on $G$, that is, linear functionals $\mathcal C_c^\infty(G)\to \mathbb C$ whose restriction to $J\times J$-invariants, for every compact open subgroup $G$, is represented by an element in the $J$-Hecke algebra $\mathcal H(G,J)$. 
\end{definition}

In other words, 
\begin{equation}
   \label{equation-completed-Hecke}
	\widehat{\mathcal H(G)} = \underset{\leftarrow}{\lim} \mathcal H(G,J),
\end{equation}
where the limit is taken over a basis of open compact subgroups, and the transition maps are given by averaging. The completed Hecke algebra is unital, with identity element $\delta_1=$ the delta measure at the identity.


\begin{proposition}
	\label{proposition-center-Hecke}
	The Bernstein center $\mathfrak Z(G)$ is naturally identified with the ring of linear endomorphisms of $\mathcal H(G)$ which commute with right and left multiplication; also, with the center of the completed Hecke algebra $\widehat{\mathcal H(G)}$.
\end{proposition}

The ``natural'' identification is the one induced by viewing $\mathcal H(G)$ as a left $G$-module (equivalently, a right $G$-module; ), so the map from $\mathfrak Z(G)$ to 


\begin{proof}
	It is easy to see that the functor $V\mapsto V^J = e_J V$ (where $e_J$ is the identity element in $\mathcal H(G,J)$) from $\mathcal M(G)$ to $\mathcal H(G,J)$-modules induces a homomorphism from the center of the former category to that of the latter. But $\mathcal H(G,J)$ is a unital algebra, so, by Example \ref{example-center-modulecategory}, $\mathfrak Z(G)$ lands in its center. The center commutes with projections (averaging) $V^{J'}\to V^J$ for $J'\subset J$, so the collection of these maps defines a homomorphism from $\mathfrak Z(G)$ to linear endomorphisms of $\mathcal H(G)$ which commute with left and right multiplication. Taking the inverse limit over $J$, we obtain similar endomorphisms of the completed Hecke algebra; but this is unital, so the ring of such endomorphisms is identified with its center. 
	
	Vice versa, every element $z$ in the center of $\widehat{\mathcal H(G)}$ projects to central elements $z_J$ in the centers of $\mathcal H(G,J)$ for all $J$, and induces functorial endomorphisms on all smooth representations $V$, by letting it act on $V^J$ as $z_J$. 
\end{proof}

The block decomposition of the category $\mathcal M(G)$ (Theorem \ref{theorem-Bernstein-decomposition}) induces a decomposition of its center,
\begin{equation}
	\label{equation-decomposition-Bernstein-center}
	\mathfrak Z(G) = \prod_\Omega \mathfrak Z(G)_\Omega,
\end{equation}
where $\Omega$ runs over all the Bernstein components, and $\mathfrak Z(G)_\Omega$ is the center of the corresponding Bernstein block $\mathfrak M(G)_\Omega$. Recall that every such component is the associate class of a pair $(L,D)$, with $L$ a Levi of $G$ and $D$ a cuspidal component of $L$, and that such a component is, as a variety, isomorphic to $D\sslash W(L,D)$.

\begin{theorem}
   \label{theorem-Bernstein-center}	
The center of the Bernstein block $\mathcal M(G)_\Omega$ is isomorphic to the ring $\mathbb C[\Omega]$ of regular functions on $\Omega$. The isomorphism is uniquely determined by the requirement that, if $\Omega$ is represented by a cuspidal component $D$ of a Levi subgroup $L$, and $P$ is a parabolic with Levi $L$, an element $z$ in the center acts on the normalized parabolic induction (Definition \ref{definition-normalized-parabolic-induction-restriction}) of a supercuspidal $\sigma\in D$ by evaluation of the corresponding regular function at the image of $\sigma$ in $\Omega = D\sslash W(L,D)$.
\end{theorem}

\begin{proof}
We follow the proof of the original paper \cite{Bernstein-centre}; a slightly different proof is contained in Bernstein's notes \cite{Bernstein-padic}, using the (imporant but heavier) machinery of {\it second adjointness}. We use normalized parabolic induction and restriction throughout. 

Firstly, the case $L=G$ was treated in Proposition \ref{proposition-endomorphism-cuspidal}.

Secondly, by Proposition \ref{proposition-cuspidal-data-equivalences}, every representation in $\mathcal M(G)_\Omega$ embeds into a sum of representations of the form $I_P(\sigma)$, for $P$ varying over parabolics with Levi $L$ and $\sigma\in \mathcal M(L)_D$, so an element of the center is determined by its action on representations of the form $I_P(\sigma)$. 

Thirdly, we can replace an arbitrary $\sigma$ with a progenerator $\Pi_D$ of the category $\mathcal M(L)_D$, for example, $	\Pi_D= \rm{cInd}_{M_0}^M \rm{Res}_{M_0}^M \pi$, as in Equation \ref{equation-progenerator-cuspidal}. We have an action of the group algebra $\mathbb C[M/M_0]=\mathbb C[\Psi_M]$ on $I_P^G(\Pi_D)$ by $G$-automorphisms (induced from its action on $\Pi_D$), and by definition the action of the Bernstein center has to commute with it. For every compact open subgroup $K$ of $G$, the space $I_P^G(\Pi_D)^K$ is a projective $\mathbb C[\Psi_M]$-module of finite rank; the action of any $z\in \mathfrak Z(G)$ on it must be given by a section of its sheaf of $\mathbb C[\Psi_M]$-linear endomorphisms. On the other hand, by Proposition \ref{proposition-irreducible-induction} for a Zariski open set of $\eta\in\Psi_G$, the specialization at $\eta$, i.e., the representation $I_P^G(\pi\otimes\eta)$, is irreducible. (Note that a supercuspidal representation with unitary central character is always unitarizable, as we can form an inviant inner product by integrating matrix coefficients on $G/Z(G)$.) It follows from Schur's lemma that $z$ must act by a scalar on this quotient. These facts combined imply that $z$ acts as an element of $\mathbb C[\Psi_M]$. Finally, if $\pi$ and $\pi\otimes\eta$ are isomorphic, then a choice of isomorphism induces an $L$-endomorphism of $\Pi_D$ (as in Proposition \ref{proposition-endomorphism-cuspidal}), hence this element should be invariant under the stabilizer $\Gamma$ of $\pi$, and descend to an element of $\mathbb C[D]$.

Fourthly, for $w\in W(L,D)$ and generic $\pi\in D$ we have an isomorphism $I_P(\pi)\simeq I_P({^w\pi})$. Similarly, for $P,Q$ with Levi $L$ and generic $\pi$ we have $I_P(\pi)=I_Q(\pi)$, again by Proposition \ref{proposition-irreducible-induction}. Therefore, this element of $\mathbb C[D]$ is independent of $P$, and invariant under $W(L,D)$. 

We have proven that we have an injection $\mathfrak Z(\mathcal M(G)_\Omega)\hookrightarrow \mathbb C[\Omega]$. There remains to prove that it is surjective. Since, again by Proposition \ref{proposition-cuspidal-data-equivalences}, we have a functorial embedding
\begin{equation}
	\label{equation-functorial-embedding-induction} 
	V\hookrightarrow \bigoplus_P I_P (R_P V)_D,
\end{equation}
and any $z\in \mathbb C[\Omega] = \mathfrak Z(\mathcal M(L)_D)^{W(L,D)}$ acts by $G$-automorphisms on the right hand side (since it acts by $L$-automorphisms on the inducing data), it is enough to show that the embedding \eqref{equation-functorial-embedding-induction} is invariant under the $z$-action. The class of representations for which this happens is stable under subquotients and direct sums, and therefore it is enough to assume that $V=I_Q(\Pi_D)$, with $\Pi_D$ as before and $Q$ with Levi $L$. More precisely, we will prove that under the canonical maps 
$$ I_Q(\Pi_D) \to I_P (R_P I_Q(\Pi_D))_D$$
the $Q$-induced action of $z$ on the left coincides with the $P$-induced action on the right. Equivalently, considering the $L$-module 
$$ \sigma= (R_P I_Q(\Pi_D))_D \in \mathcal M(L)_D,$$
the $Q$-induced action of $z$ (call it $z_Q$) coincides with its action on $\sigma$ as an object of $\mathcal M(L)_D$ (call it $z_P$). 

The geometric lemma (Proposition \ref{proposition-geometric-lemma}) provides a filtration of $\sigma$ indexed by $w\in W(L,D)$, with $w$-graded piece ${^w\Pi_D}$. The original $\mathbb C[D]$-action on the inducing datum $\Pi_D$ becomes $w$-twisted on the $w$-graded piece. 
Since $z$ is a $W(L,D)$-invariant function on $D$, the actions $z_P$ and $z_Q$ coincide on all the gradings. There remains to show that there is no unipotent difference between the two actions. Any such nontrivial unipotent difference would give rise to an $L$-equivariant morphism between two of the graded pieces, ${^{w_1}\Pi_D}$ and ${^{w_2}\Pi_D}$. Moreover, it would be $\mathbb C[D]$-equivariant with respect to the $w_1$-twisted action on the former and the $w_2$-twisted action on the latter. Again, for every compact open $K\subset L$, the $K$-invariants of ${^{w_1}\Pi_D}$ and ${^{w_2}\Pi_D}$ are finitely generated projective $\mathbb C[D]$-modules, and therefore it suffices to show that for a Zariski dense set of specializations we cannot have such morphisms. But this is indeed the case -- for a Zariski dense subset of $\pi\in D$ we have  $\Hom_L({^{w_1}\pi}, {^{w_2}\pi})=0$. It follows that $z_P$ and $z_Q$ have to coincide, completing the proof of the theorem.


\end{proof}

\section{The Satake isomorphism}
\label{section-Satake-isomorphism}


\begin{definition}
 \label{definition-unramified-group}
A reductive group over a local non-Archimedean field $F$ is said to be {\it unramified} if it is quasisplit, and splits over an unramified extension.
\end{definition}


\begin{proposition}
 \label{proposition-unramified-group}
For a connected reductive group $G$ over $F$, the following are equivalent:
\begin{enumerate}
 \item $G$ is unramified (Definition \ref{definition-unramified-group});
 \item $G$ admits a reductive model of the ring of integers $\mathfrak o$ (i.e., a smooth model with connected reductive geometric fibers).
\end{enumerate}
Moreover, the integral model over $\mathfrak o$ is unique up to $G_{\text{ad}}(F)$-conjugacy, that is, for any two reductive $\mathfrak o$-groups $\mathcal G_1$, $\mathcal G_2$ with general fiber identified with $G$, there is an isomorphism $\mathcal G_1 \simeq \mathcal G_2$ that restricts to an inner automorphism (over $F$) on $G$.
\end{proposition}

[This proposition should be moved to the chapter on algebraic groups.]

\begin{proof}
 For the direction from the second to the first, see \cite[Corollary 5.2.14]{Conrad-RGS}. The opposite direction follows from the classification in terms of root data with Galois actions. [To be added.] For the uniqueness, see \cite[Theorem 7.2.16]{Conrad-RGS}.
\end{proof}


\begin{definition}
 \label{definition-hyperspecial}
A \text{hyperspecial subgroup} of $G(F)$, where $G$ is an unramified connected reductive group over $F$, is a subgroup of the form $K = \mathcal G (\mathfrak o)$, where $\mathcal G$ is a reductive integral model.
\end{definition}

Hyperspecial subgroups are unique, up to conjugacy, for adjoint groups, as follows from the uniqueness statement of Proposition \ref{proposition-unramified-group}. This does not need to be true when $G(F)$ does not surject onto $G_{\text{ad}}(F)$.

\begin{proposition}
 \label{proposition-hyperspecial-maximal}
A hyperspecial subgroup (Definition \ref{definition-hyperspecial}) is maximal.
\end{proposition}

\begin{proof}
 
 [Omitted for now.]
\end{proof}


From now on, $G$ will denote $G(F)$. 
Fix a hyperspecial subgroup $K = \mathcal G(\mathfrak o)$, corresponding to an $\mathfrak o$-model $\mathcal G$, and consider the \emph{integral} unramified (``spherical'') Hecke algebra $\mathcal H(G,K)$ of $\mathbb Z$-valued, $K$-biinvariant functions on $G$. If we consider them as functions on the discrete space $G/K$, using the counting measure on this space we can identify them as measures, and this defines their convolution and, more generally, their action on the $K$-invariant vectors of any representation $V$ (with arbitrary coefficients!). Explicitly, if $v\in V^K$, the characteristic function of a double coset $KgK$ acts as 
$$ 1_{KgK}\cdot v = \sum_{\gamma \in [KgK/K]} \gamma\cdot v.$$

The goal of this section is to establish the integral Satake isomorphism. For this purpose, let $\mathcal Y$ be ``the'' full pre-flag variety of $\mathcal G$ over $\mathfrak o$, $\mathcal Y \simeq \mathcal N\backslash \mathcal G$, where $\mathcal N$ is the unipotent radical of a Borel subgroup. We do not really choose a Borel subgroup, but the choice of integral model matters, as it endows $\mathcal Y$ with a distinguished $K$-orbit, equal to $\mathcal Y(\mathfrak o)$, that will serve as our base point. As for the group, will use $Y$ etc.\ to denote $F$-points. We let $T\supset T_0$ denote the universal Cartan $T=B/N$, and its maximal compact subgroup $T_0=\mathcal T(\mathfrak o)$; we reserve the letter $A$ for the maximal split torus in $T$. 

We let $\mathcal S(Y/K)$ denote the space of \emph{$\mathbb Z$-valued}, compactly supported, $K$-invariant functions on $Y$. It is a module for $\mathcal H(G,K)$ (under the right action of $G$ on $Y$) and for $T$ (under the ``left'' action of $T$ on $Y$). To be clear, the action of an element $t\in T$ on functions is defined as translation by $t$, not $t^{-1}$, and it is not normalized by any modular character --- which is not defined over $\mathbb Z$: $(t\cdot f)(y) = f(ty)$; this way, the center of $G$ acts the same, whether it is considered as a subgroup of $G$ or of $T$.


\begin{lemma}
 \label{lemma-torus-module}
 Every element of $\mathcal S(Y/K)$ is $T_0$-invariant, hence the action of $T$ factors through the quotient $T/T_0$; in particular, we have an action of the Hecke algebra $\mathcal H(T,T_0)$. Under this action, $\mathcal S(Y/K)$ is a free module of rank one, and the element $e_0 = 1_{\mathcal Y(\mathfrak o)}$ is a generator.
\end{lemma}

\begin{proof}
 The Iwasawa decomposition [needs to be added] $G = N T K$ shows that $G = \bigsqcup_{t\in T/T_0} N t K$, and in particular: 
 \begin{itemize}
  \item every $N\backslash G/K$-coset is left invariant by $T_0$;
  \item the group $T/T_0$ acts simply transitively on the cosets.
\end{itemize}
One of these cosets is equal (modulo $N$) to $\mathcal Y(\mathfrak o)$.
\end{proof}

Let $\Lambda = T/T_0$.  We start with the Satake isomorphism for tori:

\begin{proposition}
 \label{proposition-Satake-torus}
Let $T$ be an unramified torus over $F$, and let $A\subset T$ be the maximal split subtorus. If $\varpi \in F$ is a uniformizer, the map $\Lambda:=X_*(A)\ni \lambda\mapsto \lambda(\varpi) \in A(F)$ descends to an isomorphism $X_*(A)\simeq T/T_0$. 

Moreover, let $\check T$ be the dual torus to $T$, understood as a group scheme over $\mathbb Z$, with an action of the unramified Galois group $\Gamma = \left< \sigma \right>$, where $\sigma$ denotes the Frobenius element. Then, the dual $\check A$ of $A$ is the maximal torus quotient of $\check T$ where $\Gamma$ acts trivially, and the natural maps induce isomorphisms of algebras
\begin{equation}
 \label{equation-Satake-torus}
 \mathbb Z[\check T \sigma]^{\check T} = \mathbb Z[\check A] = \mathbb Z[\Lambda] = \mathcal H(T,T_0),
\end{equation}
where the ``coset'' $\check T \sigma$ is the space $\check T$ equipped with the $\sigma$-twisted conjugation of $\check T$, $x\sigma \cdot t = (t\cdot {^\sigma t^{-1}}) x \sigma$
(and the notation $\mathbb Z[\cdot]$ is used both for group rings and coordinate rings --- it should be clear which is which).
\end{proposition}

\begin{proof}
 Consider the quotient of algebraic group schemes over $\mathfrak o$:
 $$ 1\to \mathcal A \to \mathcal T \to \mathcal T/\mathcal A\to 1.$$
 Since $A$ is split, this induces (by Hilbert 90) surjections at the level of $F$- and $\mathbb F_q$-points, hence (by smoothness) at the level of $\mathfrak o$-points.  On the other hand, $T/A$ is anisotropic, and $T/A(F)=\mathcal T/\mathcal A(\mathfrak o)$. Therefore, $T(F) = A(F) \mathcal T(\mathfrak o)$, and since $A(F)\cap \mathcal T(\mathfrak o) = A(\mathfrak o)$, and $A(F)/A(\mathfrak o) = \Lambda$, this shows the bijection $\Lambda \xrightarrow\sim T/T_0$.
 
 For the coordinate rings, if $\Lambda'$ is the $\bar F$-cocharacter group of $T$, then the torus $A$ is spanned by the images of Galois-stable cocharacters, hence $\Lambda = (\Lambda')^\Gamma$. On the dual side, the embedding $\Lambda \hookrightarrow \Lambda'$ induces a morphism of dual tori $\check T \to \check A$, which identifies $\check A$ with the quotient of $\check T$ by the subtorus of all elements of the form $(t\cdot {^\sigma t^{-1}})$, $t\in \check T$, hence   
 $\mathbb Z[\check T\sigma ]^{\check T} = \mathbb Z[\check A] = \mathbb Z[\Lambda]$.
\end{proof}

Now denote by $t_\lambda$ a representative for $\lambda \in \Lambda = T/T_0$ in $T$, and let $\delta_\lambda=1_{N t_\lambda K} = t_\lambda^{-1}\cdot 1_{\mathcal Y(\mathfrak o)}$; thus, the elements $\delta_\lambda$ form a basis for $\mathcal S(Y/K)$. By Lemma \ref{lemma-torus-module}, the action map $\mathcal H(T,T_0)\ni h \mapsto h\cdot 1_{\mathcal Y(\mathfrak o)}$ identifies the spaces $\mathcal S(Y/K)$ and $\mathcal H(T,T_0)\simeq \mathbb Z[\Lambda]$; notice, however, that the characteristic function $t_\lambda T_0$ in $\mathcal H(T,T_0)$ corresponds to $\delta_{-\lambda}$. 

\begin{theorem}[Satake isomorphism]
\label{theorem-Satake-isomorphism}
 Let $S: \mathcal H(G,K) \to \mathcal H(T,T_0)$ be the map given by the action map $\mathcal H(G,K)\ni h \mapsto h\cdot 1_{\mathcal Y(\mathfrak o)} \in \mathcal S(Y/K)$ and the identification of $\mathcal S(Y/K)$ with $\mathcal H(T,T_0)$ (again through the analogous action map), i.e., $S(h)$ is that element of $\mathcal H(T,T_0)$ such that $h\cdot 1_{\mathcal Y(\mathfrak o)} = S(h)\cdot 1_{\mathcal Y(\mathfrak o)} $. Then $S$ gives rise to an isomorphism of algebras: 
 \begin{equation}
 \label{equation-Satake}
 \mathcal H(G,K) \simeq \mathcal H(T,T_0)_{\mathbb Z} \cap \mathcal H(T,T_0)_{\mathbb Q}^{W,\bullet} = \mathbb Z[\Lambda]\cap \mathbb Q[\Lambda]^{W,\bullet},
 \end{equation}
 where the $\bullet$-action of the (relative) Weyl group $W$ is defined by $(w\bullet f)(t_\lambda)  = q^{\left < \lambda, \rho - w\rho \right>} f(t_{w^{-1}\lambda})$ (where $q$ is the degree of the residue field).
\end{theorem}

Explicitly, the element $h \in \mathcal H(G,K)$, is mapped to the element of $\mathcal H(T,T_0)$ whose value at $t$ is equal to 
\begin{equation}
 \label{equation-Satake-integral}
h\cdot 1_{\mathcal Y(\mathfrak o)} (Nt^{-1}) = \int_G h(g) 1_{NK} (Nt^{-1}g) dg =  \int_N h(tn) dn,
\end{equation}
where the Haar measure on $N$ gives volume $1$ to $\mathcal N(\mathfrak o)$. The way that the Satake isomorphism is usually defined in the literature is through the equation \eqref{equation-Satake-integral}, multiplied by $\delta_B(t)^{\frac{1}{2}}$ (where $\delta_B$ is the modular character of the Borel subgroup), in order to replace the $\bullet$-action of $W$ by the usual action of $W$ --- but this modification is not defined over $\mathbb Z$.

Notice also that, the characteristic function of $t_\lambda T_0$ in $\mathcal H(T,T_0)$ corresponds to the element $\delta_{-\lambda}$ of $\mathcal S(Y/K)$, image of the action map will be invariant under the following Weyl group action on $\mathcal S(Y/K)_{\mathbb Q}$:
\begin{equation}
 \label{equation-bullet-Weyl}
w\bullet \delta_{\lambda} = q^{\left < \lambda, \rho - w^{-1}\rho\right>} \delta_{w\lambda}.
\end{equation}

\begin{proof}
 First of all, we notice that the map $S$ is a homomorphism of algebras, because the actions of $G$ and $T$ on $Y$ commute, and $\mathcal H(T,T_0)$ is abelian: if $h_i\in \mathcal H(G,K)$ and $h_i'\in \mathcal H(T,T_0)$ are such that $h_i\cdot  1_{\mathcal Y(\mathfrak o)} = h_i' \cdot 1_{\mathcal Y(\mathfrak o)}$, then $(h_1 h_2) \cdot 1_{\mathcal Y(\mathfrak o)} = h_1 \cdot (h_2 \cdot 1_{\mathcal Y(\mathfrak o)}) = h_1 ( h_2' \cdot 1_{\mathcal Y(\mathfrak o)}) = h_2' (h_1\cdot 1_{\mathcal Y(\mathfrak o)}) = (h_2'  h_1') \cdot 1_{\mathcal Y(\mathfrak o)} = (h_1'h_2')\cdot 1_{\mathcal Y(\mathfrak o)}$.
 
 Next, we claim that the image lies in $\mathbb Q[\Lambda]^{W,\bullet}$. We will present two proofs for that, one of them only for the split case [but it can be generalized --- to do]. 
 
 The first proof, following \cite{Cartier-Corvallis}, uses the explicit expression \eqref{equation-Satake-integral} of the Satake transform as an integral, and replaces it by an orbital integral for the action of $G$ on itself by conjugacy. Namely, choose a lift of the quotient map $B\to T$, thus identifying $T$ as a subtorus of $B$. If $t \in T$ is a \emph{regular} element (i.e., has trivial stabilizer under the action of the Weyl group), then we have the following formula:
 $$ \int_N h(t n) dn =  |\det(\text{Ad}_{\mathfrak n} (t^{-1}) - 1)| \int_N f(n^{-1} t n) dn = |\det(\text{Ad}_{\mathfrak n} (t^{-1}) - 1)| \int_{T\backslash G} f(g^{-1} t g) dg,$$
 where $\text{Ad}_{\mathfrak n}$ denotes the left adjoint action of $T$ on the Lie algebra $\mathfrak n$, and the invariant measure on $T\backslash G$ is normalized so that the total measure of $K$-orbits represented by $T\backslash TN(\mathfrak o)$ is $1$. This formula follows easily by considering the map $N\to tN$ given by $n\mapsto n^{-1} t n$, and representing the measures as absolute values of volume forms; at the last step, one uses the $K$-invariance of $H$ to represent $T\backslash G$ by $NK$. 
 
 Hence, for $w\in W$ and $t$ a \emph{regular} element in $T$ (any class in $T/T_0$ has such representatives),
 $$ \frac{\int_N h(n{^wt}) dn}{\int_N h(nt) dn} = \frac{|\det(\text{Ad}_{\mathfrak n} ({^wt^{-1}}) - 1)|}{|\det(\text{Ad}_{\mathfrak n} (t^{-1}) - 1)|} = \left | \prod_{\alpha > 0}  \frac{1-e^{-w^{-1}\alpha}(t)}{1-e^{-\alpha}(t)} \right |$$
 $$ = \left| \prod_{\alpha > 0, w\alpha <0}  \frac{1-e^{\alpha}(t)}{1-e^{-\alpha}(t)} \right| = \left| \prod_{\alpha > 0, w\alpha <0} e^{\alpha}(t)\right| = | e^{\rho -w\rho}(t)|,$$
 which amounts to the stated invariance property. (We have not assumed $G$ to be split, for this calculation: the terms inside the absolute values are algebraic functions, and therefore it is valid to manipulate them over the algebraic closure --- where all roots are defined.)
 
 [Another proof with Fourier transforms to be added.]
 
 Finally, we prove that the map $\mathcal H(G,K) \to  \mathbb Z[\Lambda]\cap \mathbb Q[\Lambda]^{W,\bullet}$ is an isomorphism. We will argue by identifying the space on the right (call it $M$) as a subspace of $\mathcal S(Y/K)$ --- the $\bullet$-action of $W$ is given by \eqref{equation-bullet-Weyl}. Notice that $M$ is a free $\mathbb Z$-module with generators $m_\lambda$, indexed by dominant cocharacters $\lambda$ into $A$, given by $m_\lambda= \sum_{\lambda'\sim \lambda} q_{\lambda'} \delta_{\lambda'}$, where $\lambda' \sim \lambda$ means that $\lambda'=w\lambda$ for some $w\in W$, and in this case we set $q_{\lambda'} = q^{\left < \lambda, \rho - w^{-1}\rho\right>}$. We define a filtration of this module with respect to the partial ordering by coroots, $\mu \ge \lambda$ if $\mu-\lambda$ is a sum of positive coroots. Similarly, $\mathcal H(G,K)$ has a basis consisting of the characteristic functions of the cosets $Kt_{-\lambda} K$ (with $\lambda$ dominant, again), and we use it to define a filtration of $\mathcal H(G,K)$ indexed by dominant weights. We claim that the map $\mathcal H(G,K)\to M$ respects these filtrations:
 $$ F^\lambda \mathcal H(G,K)\to F^\lambda M,$$
 and that the generator of the $\lambda$-th graded piece, represented by the function $1_{Kt_\lambda K}$, maps to the generator of the $\lambda$-th graded piece, represented by $m_\lambda$. These statements follow from the following fundamental fact:
 
 \begin{quote}
  For $\lambda$ dominant, we have 
\begin{equation}
\label{equation-Hecke-action-preflag} 
K t_\lambda K \subset \bigcup_{\mu\le \lambda} N t_\mu\cdot K,
\end{equation}
and $K t_\lambda K \cap N t_\lambda K = \mathcal N(\mathfrak o)  t_\lambda  K$. 
 \end{quote}

 [The proof of this will be added together with the proof of the Cartan and Iwasawa decompositions.]
 
This implies that $1_{K t_{-\lambda} K} \cdot \delta_0 = \delta_\lambda + \sum_{\mu<\lambda} c_{\mu,\lambda} \delta_\mu$ for some coefficients $c_{\mu, \lambda}\in\mathbb N$. We leave it to the reader to check that this is equivalent to the claim.
 
\end{proof}

\section{Langlands parameters}
\label{section-Langlands-parameters}

Let $G$ be a (connected) reductive group over a local field $F$. We will write $G$ for $G(F)$. The $L$-group and the $C$-group of $G$ have been defined in Section \ref{reductiveforms-section-Lgroups}. We denote by $\Gamma_F$ the Galois group of $F$ (of a fixed separable extension\footnote{See Remark \ref{reductiveforms-remark-Lgroup-sheaf}: it is better not to fix a separable extension, and to translate these definitions to sheaves over the \'etale site of $F$.}), and by $\mathcal W_F$ its Weil group. For definitions, see \cite{Tate-Corvallis}. We only remind here that the Weil group comes with isomorphisms $\mathcal W_F/\mathcal W_E = \Gamma_F/\Gamma_E = \Hom(E,F^s)$ for every separable extension $E$ of $F$, and $\mathcal W_F^{ab}\xrightarrow\sim F^\times$, compatible with the isomorphism $\Gamma_F \xrightarrow\sim \widehat{F^\times}$ (profinite completion) of class field theory. As in \cite{Tate-Corvallis}, we will normalize the isomorphism of class field theory so that a Frobenius element maps to the \text{inverse} of a uniformizer, i.e., a \emph{geometric Frobenius} element maps to a uniformizer. In particular, we have a norm map $|\bullet|: \mathcal W_F \to  F^\times \to \mathbb R_+^\times$, sending a Frobenius element to $q$: the degree of the residue field.

We also remind of the modification of the Weil group that is needed in order to pass from $l$-adic to complex representations:

\begin{definition}
 \label{definition-Weil-Deligne}
Let $F$ be a non-Archimedean field. The {\it Weil--Deligne group} $\mathcal W_F'$ is the semidirect product $\mathcal W_F\ltimes \mathbb G_a$, with $w x w^{-1} = |w| x$ for $w\in \mathcal W_F$ and $x\in \mathbb G_a$. A {\it representation of the Weil--Deligne group} over a field $E$ of characteristic zero is a pair $(\rho, N)$ consisting of a representation of $\mathcal W_F$ \emph{with open kernel} on a finite-dimensional vector space $V$ over $E$, and a nilpotent endomorphism $N$ of $V$, satisfying $\rho(w)N\rho(w)^{-1} = |w|\dot N$. 
\end{definition}

The ``open kernel'' condition is the important one here; it makes irrelevant the topology of $\text{GL}_E(V)$. The Weil--Deligne group is a convenient way to de-topologize the $l$-adic representations of the Weil group that show up ``in nature'' (in \'etale cohomology), and translate them among different $l$'s, or to the complex numbers:

\begin{proposition}
 \label{proposition-Weil-Deligne-representations}
Let $l$ be a prime different from the residual characteristic $p$ of (a non-Archimedean field) $F$, and let $E$ be a finite extension of $\mathbb Q_l$. There is a canonical bijection between isomorphism classes of (continuous) finite-dimensional $E$-representations $\phi:\mathcal W_F \to \text{GL}(V)$ and representations $(\rho, N)$ of the Weil--Deligne group over $E$ (Definition \ref{definition-Weil-Deligne}), characterized by the property that 
\begin{equation}
 \label{equation-Weil-Deligne-representations}
\phi(\Phi\sigma) = \rho(\Phi\sigma) \exp(t_l(\sigma) N),
\end{equation}
for some Frobenius element $\Phi\in \mathcal W_F$, any element $\sigma$ of the inertia subgroup, and $t_l$ a choice of isomorphism of the pro-$l$-quotient of (tame) inertia with $\mathbb Z_l$.
\end{proposition}

Recall that the tame inertia quotient is generated by $n$-th roots of a uniformizer, for $(n,p)=1$, and is isomorphic (up to a choice of topological generator) to $\widehat{\mathbb Z}^p = \prod_{l\ne p} \mathbb Z_l$.

\begin{proof}
 See \cite[\S 4.2]{Tate-Corvallis} for references.
\end{proof}






\begin{definition}
 \label{definition-Langlands-parameter}
A {\it Langlands parameter} into the $L$-group of $G$ is a morphism $\mathcal W_F'\to{^LG}$ over $\Gamma$.
\end{definition}




The local Langlands conjecture posits the existence of a canonical finite-to-one map:

$$\{\mbox{irreducible admissible representations of }G\}/\sim \to \{\mbox{Langlands parameters into }{^LG}\}/\sim.$$


%**************************************************************************************


%*************************************************************************




\begin{multicols}{2}[\section{Other chapters}]
\noindent
\begin{enumerate}
\item \hyperref[introduction-section-phantom]{Introduction}
\item \hyperref[representationtheory-section-phantom]{Basic Representation Theory}
\item \hyperref[representations-compact-section-phantom]{Representations of compact groups}
\item \hyperref[liegroups-general-section-phantom]{Lie groups and Lie algebras: general properties}
\item \hyperref[liestructure-section-phantom]{Structure of finite-dimensional Lie algebras}
\item \hyperref[algebraicgroups-section-phantom]{Linear algebraic groups}
\item \hyperref[reductiveforms-section-phantom]{Forms and covers of reductive groups, and the $L$-group}
\item \hyperref[vermamodules-section-phantom]{Verma modules}
\item \hyperref[representations-local-section-phantom]{Representations of reductive groups over local fields}
%\item \hyperref[gKmodules-section-phantom]{$(\mathfrak g, K)$-modules}
%\item \hyperref[asymptotics-section-phantom]{Asymptotics and the Langlands classification}
\item \hyperref[plancherel-section-phantom]{Plancherel formula: reduction to discrete spectra}
\item \hyperref[discreteseries-section-phantom]{Construction of discrete series}
\item \hyperref[galoiscohomology-section-phantom]{Galois cohomology of linear algebraic groups}
\item \hyperref[automorphicspace-section-phantom]{The automorphic space}
%\item \hyperref[harmonicanalysis-section-phantom]{Harmonic analysis over local fields}
%\item \hyperref[automorphicforms-section-phantom]{Automorphic forms}
%\item \hyperref[periods-section-phantom]{Periods, theta correspondence, related methods}
%\item \hyperref[traceformulalocal-section-phantom]{The trace formula: local aspects}
%\item \hyperref[traceformulaglobal-section-phantom]{The trace formula: global aspects}
%\item \hyperref[arithmetic-section-phantom]{Arithmetic, reciprocity, Shimura varieties}
%\item \hyperref[geometric-section-phantom]{Geometric aspects}
\item \hyperref[fdl-section-phantom]{GNU Free Documentation License}
\item \hyperref[index-section-phantom]{Auto Generated Index}
\end{enumerate}
\end{multicols}





\bibliography{my}
\bibliographystyle{amsalpha}

\end{document}

